% Options for packages loaded elsewhere
\PassOptionsToPackage{unicode}{hyperref}
\PassOptionsToPackage{hyphens}{url}
%
\documentclass[
  pdflatex,
  sn-mathphys-num]{sn-jnl}
\usepackage{amsmath,amssymb}
\usepackage{iftex}
\ifPDFTeX
  \usepackage[T1]{fontenc}
  \usepackage[utf8]{inputenc}
  \usepackage{textcomp} % provide euro and other symbols
\else % if luatex or xetex
  \usepackage{unicode-math} % this also loads fontspec
  \defaultfontfeatures{Scale=MatchLowercase}
  \defaultfontfeatures[\rmfamily]{Ligatures=TeX,Scale=1}
\fi
\usepackage{lmodern}
\ifPDFTeX\else
  % xetex/luatex font selection
\fi
% Use upquote if available, for straight quotes in verbatim environments
\IfFileExists{upquote.sty}{\usepackage{upquote}}{}
\IfFileExists{microtype.sty}{% use microtype if available
  \usepackage[]{microtype}
  \UseMicrotypeSet[protrusion]{basicmath} % disable protrusion for tt fonts
}{}
\makeatletter
\@ifundefined{KOMAClassName}{% if non-KOMA class
  \IfFileExists{parskip.sty}{%
    \usepackage{parskip}
  }{% else
    \setlength{\parindent}{0pt}
    \setlength{\parskip}{6pt plus 2pt minus 1pt}}
}{% if KOMA class
  \KOMAoptions{parskip=half}}
\makeatother
\usepackage{xcolor}
\usepackage{longtable,booktabs,array}
\usepackage{calc} % for calculating minipage widths
% Correct order of tables after \paragraph or \subparagraph
\usepackage{etoolbox}
\makeatletter
\patchcmd\longtable{\par}{\if@noskipsec\mbox{}\fi\par}{}{}
\makeatother
% Allow footnotes in longtable head/foot
\IfFileExists{footnotehyper.sty}{\usepackage{footnotehyper}}{\usepackage{footnote}}
\makesavenoteenv{longtable}
\setlength{\emergencystretch}{3em} % prevent overfull lines
\providecommand{\tightlist}{%
  \setlength{\itemsep}{0pt}\setlength{\parskip}{0pt}}
\setcounter{secnumdepth}{-\maxdimen} % remove section numbering
\ifLuaTeX
  \usepackage{selnolig}  % disable illegal ligatures
\fi
\usepackage[]{natbib}
\usepackage{bookmark}
\IfFileExists{xurl.sty}{\usepackage{xurl}}{} % add URL line breaks if available
\urlstyle{same}
\hypersetup{
  hidelinks,
  pdfcreator={LaTeX via pandoc}}

\author{}
\date{}

\begin{document}
\title{Limits of Independent Security Verification on Vendor-Locked Devices}
\author{\fnm{Patrick} \sur{Schleizer}}
\abstract{Vendor-locked devices -\/- most prominently stock Android phones, which
come with a locked bootloader and no administrative rights -\/- are
general-purpose computers whose owners cannot fully inspect, control, or
verify the software they run. This paper introduces the concept of
owner-inspectability: the owner\textquotesingle s ability to examine the
whole system from a clean, trusted vantage point and demonstrate what it
runs to a verifier of their choosing. It adopts the
owner\textquotesingle s perspective, treating security as the ability to
audit the device rather than the platform\textquotesingle s ability to
enforce policy. Today, several limitations prevent this. The owner
cannot inspect the system from a clean environment or directly access
its storage and memory. Application and update delivery is
vendor-controlled, so compromise or coercion could enable targeted
malicious installations. Hardware-backed attestation lets apps and
services check that the device reports an unmodified, vendor-approved
state. Modifying it for inspection therefore makes the device fail
attestation, causing banking, payment, and similar apps that depend on
it to refuse to run. Root access does not fully remove these
limitations: it provides only an in-system vantage point, not control
over the boot chain and underlying firmware, and can be deceived by a
kernel-space rootkit. The owner is caught in a bind: inspecting the
device requires the very modification that attestation rejects. Most
established malware-analysis techniques are thus unavailable without
first breaking the phone\textquotesingle s attestation guarantees,
making the absence of findings weak evidence of integrity. The paper
asks what capability would escape this bind: owner-controlled privileged
inspection that preserves verified boot and attestation. It finds that
no shipped or proposed approach provides this capability, leaving an
open problem. The analysis focuses on stock Android but extends, to
varying degrees, to other vendor-locked devices.}
\keywords{Android security, independent verification, mobile forensics, remote attestation, rootkit detection, malware analysis, device integrity}
\maketitle

\subsection{1. Introduction}\label{introduction}

On a vendor-locked device the owner is, in practice, obliged to trust
the platform operator and device vendor, because the device cannot be
fully examined from the outside while it remains locked down. On a stock
Android phone that means trusting Google for application delivery
through the Play Store \citep{google-play-help} and the device
manufacturer for system updates. On a non-Google device such as a
Samsung those updates are the manufacturer\textquotesingle s own rather
than Google\textquotesingle s, though Google\textquotesingle s services
retain a privileged role in the update path on devices certified for
Google Mobile Services \citep{aosp-sepolicy-gmscore} (Section 6). If
that trust is abused -\/- through compromise of the delivery
infrastructure, coercion of an operator, or a malicious insider -\/-
targeted malicious installations or updates may be difficult to detect,
independently verify, or prove to a third party.

The same constraints bear on malware investigation, though their force
depends on the class of malware. Commodity, application-level threats
can often be detected on the device itself, by consumer antivirus
\citep{kaspersky-android} or Google\textquotesingle s Play Protect
\citep{google-play-protect}, without unlocking or modifying it -\/- a
point Section 12 takes up in full. This limitation is most acute for the
targeted, low-level compromise: a kernel-space rootkit or
sub-operating-system implant masks its own presence and sits beneath any
in-system scanner, which it can evade whatever that
scanner\textquotesingle s privileges (Section 9). Against that class the
owner typically cannot carry out the standard forensic approach -\/-
examining the device from a separate, known-clean computer and working
from a write-blocked image rather than the live original, as standard
forensic practice directs \citep{nist-sp800-101r1},
\citep{nist-sp800-86} -\/- without first unlocking or modifying the
device. That modification can itself alter the evidence under
examination, the mandatory wipe triggered by unlocking the bootloader
being the extreme case \citep{aosp-bootloader-locking}; and a
well-constructed backdoor could detect the onset of an investigation and
erase itself or conceal its traces \citep{chen-anti-virtualization},
\citep{ruggia-evasive-malware}, \citep{garfinkel-anti-forensics}, as
Section 14 develops.

This paper treats these properties not as implementation accidents but
as consequences of a design in which administrative capability is
deliberately withheld from the owner. It catalogs the resulting
verification limits, grounds its claims in the available technical and
forensic literature where such sources exist, and distinguishes claims
supported by cited evidence from claims that remain asserted without a
specific source.

This is a systematization-and-position paper: it introduces no new
experiment or dataset and instead organizes existing technical and
forensic evidence around a single owner-centric criterion. Its
contributions are:

\begin{itemize}
\tightlist
\item
  it names and defines owner-inspectability as the property an owner
  needs in order to audit a device, and distinguishes it from the
  related but distinct notion of organizational forensic readiness;
\item
  it catalogs, against that criterion, the specific verification limits
  of stock Android and traces each to the same deliberate design choice
  rather than to isolated implementation defects;
\item
  it shows that root access -\/- as distinct from the full
  administrative control withheld from the owner -\/- does not restore
  owner-inspectability on a modern hardened device, because root is only
  an in-system vantage point: it does not reach the kernel, firmware,
  and boot chain beneath it, it can be deceived by a kernel-space
  rootkit that falsifies what any in-system tool reports regardless of
  privilege, and reaching it forfeits the attestation standing the owner
  seeks to preserve;
\item
  it documents how remote attestation, and the applications that depend
  on it, penalize owner-controlled configurations by a mechanism that
  has tightened over time, deepening the circular trust problem;
\item
  it provides a structured capability comparison with a fully controlled
  host that makes the aggregate effect legible as a single design
  trade-off; and
\item
  it frames owner-controlled privileged inspection -\/- a verified
  owner-inspection mode that preserves verified boot and relying-party
  attestation -\/- as the capability that would be needed to escape the
  inspection-versus-attestation bind, systematizes the design space of
  approaches shipped or proposed, and identifies the gap no current
  approach closes.
\end{itemize}

The paper is organized as follows. Section 2 summarizes the platform
mechanisms. Section 3 states the threat model, scope, and terminology,
including the constrained meaning of root and the definition of
owner-inspectability. Section 4 situates the work in prior literature.
Sections 5 through 11 catalog the verification limits. Section 12
addresses objections, including aftermarket operating systems. Section
13 asks whether a verified owner-inspection mode could escape the bind,
and maps the design space of attempts. Section 14 discusses the findings
and their limitations, and Section 15 concludes.

\subsection{2. Background}\label{background}

Android\textquotesingle s security model is built around several
mutually reinforcing mechanisms. Android Verified Boot (AVB) establishes
a chain of trust from an immutable hardware root of trust through the
bootloader and operating-system partitions, so that only software signed
by the vendor executes during startup \citep{aosp-verified-boot}; a
locked bootloader refuses to load or boot any image that is not so
signed \citep{aosp-bootloader-locking}. Application sandboxing isolates
each application, giving it a distinct kernel-level user identity and a
private data directory that other applications cannot read
\citep{aosp-app-sandbox}. Mandatory Access Control, implemented through
SELinux, confines even privileged system processes to narrowly defined
domains, so that a compromise of one component does not automatically
yield access to the whole system \citep{aosp-selinux}. File-based
encryption, mandatory since Android 10, protects user data at rest, with
the principal encryption key bound to dedicated hardware
\citep{aosp-file-based-encryption}.

By design, these mechanisms constrain what ordinary software, including
most malicious software, can do once it is running on the device
\citep{mayrhofer-android-security-model}. That constraint is real but
neither symmetric nor complete. A sufficiently capable adversary is not
reliably excluded by it -\/- as Sections 8 and 12 develop, such an
adversary can compromise the already-booted, legitimately signed system
at runtime, or arrive inside the vendor-signed image itself, without the
device failing attestation, and verified boot has demonstrated defeats
beneath the operating system \citep{bootstomp},
\citep{shakevsky-trustzone}.

The same design, meanwhile, withholds from the device owner the
administrative privileges that would permit raw disk access, live memory
capture, kernel instrumentation, or the booting of an external trusted
environment \citep{aosp-adb-root}, \citep{vidas-android-collection}. The
owner is therefore placed not in the same position as an attacker but,
with respect to introspection, in a more constrained one, for a
structural rather than a measured reason: the owner\textquotesingle s
only route to inspecting the device is the modification these
protections are built to bar, whereas the capable adversary that the
owner most fears need not modify the device at all and keeps the
runtime, supply-chain, and sub-operating-system routes just described.
These controls therefore foreclose the owner\textquotesingle s only
inspection route -\/- and that of a security researcher acting on the
owner\textquotesingle s behalf -\/- while leaving the runtime,
supply-chain, and sub-operating-system routes of a capable adversary
unaffected.

These mechanisms serve the owner\textquotesingle s protection, but they
also serve platform-control and relying-party objectives that the owner
may not share. The platform works to keep unauthorized and malicious
software off the device, through a layered apparatus, each layer of
which covers a specific surface: verified boot and a locked bootloader
admit only vendor-signed code to the boot and system partitions, keeping
unauthorized code out of the system image though not out of the
application layer \citep{aosp-verified-boot},
\citep{aosp-bootloader-locking} (a chain that has itself demonstrated
sub-operating-system defeats; Section 12); Google Play Protect screens
apps and works to prevent the installation of harmful ones, with
installation from unknown sources restricted by default
\citep{google-play-protect}, \citep{nistSP800124r2}; and the Play
Integrity API gates relying-party access when a device appears tampered
\citep{play-integrity-verdicts}.

How well the app-layer screen catches malware is assessed here using
independent testing results: in AV-Comparatives\textquotesingle{} May
2024 testing Play Protect reached a 99.3 percent protection rate,
clearing the 99 percent certification threshold that round but trailing
the leading third-party scanners at 99.9 percent and logging more false
positives, an improvement over earlier independent rounds that had
ranked it lowest, at well below 70 percent detection
\citep{avcomparatives-mobile-2024},
\citep{avtest-android-endurance-2018}.

But the signed-image and attestation machinery also enforces objectives
the owner may neither share nor be able to remove or audit -\/- that
only a vendor-approved configuration runs, with application and service
access gated by relying parties (the remote apps and services that
consume the attestation verdict) rather than by the owner
\citep{play-integrity-verdicts} -\/- so the guarantee runs in part to
the platform and those relying parties, not to the owner alone. Such
functionality, which the owner may not want, cannot remove, and cannot
audit, has been termed an anti-feature in the free-software literature
\citep{hill2007antifeatures}, \citep{fdroid-antifeatures}; preinstalled
applications, advertising components, and tracking software that the
owner cannot uninstall without root are a documented instance
\citep{gamba2020preinstalled}; Carrier IQ, preinstalled carrier
diagnostic software that wrote dialer keypresses and SMS text to system
logs on some devices, is an earlier one \citep{eckersley-carrieriq}.

Content protection is the platform-control counterpart on the content
side: on devices with the requisite hardware support, hardware-backed
digital-rights-management can bind decryption and the decoded video
frames to the same hardware chain of trust as verified boot, routing
protected content through a trusted output path so that the owner may
play licensed media but cannot directly extract the decoded stream,
which is released only on terms the content provider sets
\citep{aosp-drm}. This paper treats those platform-control and
content-protection objectives as out of scope to analyze (Section 3) and
returns to the question of which party the attestation guarantee serves
(Section 8).

\subsection{3. Threat Model and Scope}\label{threat-model-and-scope}

This paper defines security from the perspective of the end user and
device owner: the ability to control, inspect, and independently verify
the device. Under this definition a device is more secure when its owner
can establish, from a trusted vantage point, whether it is running only
the software they intend and what that software is doing -\/- free of
malware or backdoors, whether planted by an outside attacker or by the
platform itself -\/- and can produce evidence of that state that a
skeptical third party could check.

Three practical capabilities are treated as the goals of interest:
inspecting the device deeply for malware and backdoors; proving the
result to a skeptical third party; and creating complete device backups
-\/- a bit-for-bit image of the whole device, all partitions and user
data, not the platform\textquotesingle s limited app-level backup -\/-
suitable for independent verification \citep{nist-sp800-101r1}. Proving
the result of that deep examination to a skeptical third party is, on a
general-purpose computer, the routine of forensic soundness: the
examiner makes a write-blocked, bit-for-bit image of the storage and
computes a cryptographic hash of it, which establishes the integrity of
the evidence -\/- that the image is byte-identical to the original -\/-
and makes the analysis reproducible on exactly those bits
\citep{nist-sp800-86}, \citep{rfc3227}. This does not by itself prove
the device is clean; what it proves is that the skeptic need not trust
the examiner\textquotesingle s verdict, because the owner holds the
published hash and a documented procedure and can hand any independent
examiner the same verified bits to re-run the malware-and-backdoor
analysis and reach their own conclusion.

Explicitly out of scope are digital-rights-management objectives and
platform-control goals -\/- that is, the platform\textquotesingle s
interest in enforcing policy that constrains the owner
\citep{doctorow-general-computation}, in guaranteeing content protection
\citep{aosp-drm}, or in ensuring that only vendor-approved
configurations run \citep{aosp-verified-boot},
\citep{aosp-bootloader-locking}. These objectives frequently stand in
tension with owner-centric verifiability, and conflating the two
obscures the specific limitations examined here. Where the platform
regards a property as a security feature and the owner regards the same
property as an obstacle to verification, this paper describes it from
the owner\textquotesingle s standpoint and notes the
platform\textquotesingle s view where relevant.

Although this paper documents Android in detail as its worked example
-\/- it is among the most openly specified of the mainstream
vendor-locked platforms, the core of its system software released as the
open-source Android Open Source Project (AOSP) \citep{aosp-open-source}
-\/- that openness is a property of the source, not of the device as
sold. AOSP is a codebase, not a complete operating system for real
hardware: running it on a phone requires device-specific proprietary
binaries the vendor distributes separately -\/- drivers, radio and
baseband firmware, and hardware-abstraction blobs under non-open
licenses \citep{google-driver-binaries}, \citep{aosp-hal} -\/- and the
mainstream "Google Android" product adds Google Mobile Services, a
proprietary suite that is not part of AOSP and available only under
license from Google \citep{google-gms}. The device as sold is therefore
vendor-locked despite the open-source codebase beneath it, bundling
proprietary components the owner cannot rebuild and shipping with the
bootloader locked \citep{aosp-bootloader-locking} (Sections 8 and 12).
The owner-inspectability criterion and the limits it exposes are not
Android-specific. They follow from vendor-lock itself: a locked boot
chain, administrative capability withheld from the owner, and
attestation that ties the device\textquotesingle s standing with apps
and services to a vendor-approved state. By the same logic they apply,
to varying degrees, to other vendor-locked devices, Apple iOS among
them; the mechanisms differ in detail, but the owner\textquotesingle s
structural position does not. This paper references such platforms only
where a cited source speaks to them and does not document them to the
same depth.

Terminology is used consistently throughout. A device is vendor-locked,
in the sense used here, when as sold it runs only vendor-signed system
images under signing keys the vendor holds and the owner does not, with
the bootloader locked -\/- so that the owner\textquotesingle s only
route to inspecting or replacing that software is to unlock the device,
which, where the vendor permits it at all, triggers a mandatory wipe of
user data \citep{aosp-bootloader-locking} and forfeits the verified-boot
and hardware-backed attestation standing the device reports to software
\citep{aosp-verified-boot} (Sections 8 and 12). The term names a
structural position -\/- locked and vendor-keyed by default, the only
exit being destructive and, for inspection, counterproductive -\/- not
the narrower claim that unlocking is impossible: a Google Pixel permits
unlocking yet is vendor-locked in this sense, while Apple iOS, which
offers owners no unlock at all, is the limiting case. Stock Android
denotes a typical device as sold, with a locked bootloader
\citep{aosp-bootloader-locking} and without root access
\citep{vidas-android-collection}.

Root access denotes uid-0 authority within the running operating system:
the ability, in principle, to read and write files, inspect processes,
and alter system components that an ordinary user cannot. The term is
inherited from Unix, where "root" is the superuser account; Android,
however, exposes no such login account. uid 0 exists as a kernel
identity held by init and a deliberately shrinking set of system daemons
\citep{aosp-root-process-isolation}, but a stock production ("user")
build ships no owner-accessible "su" and refuses to run its debug bridge
(adbd) as root \citep{aosp-adb-root}, so the owner cannot assume that
identity at all without first modifying the device.

Equally important is what this authority does not include on a modern
hardened device, because the common assumption that root confers total
control of the machine no longer holds. Root runs inside, and is bounded
by, the platform\textquotesingle s other defenses. Mandatory access
control through SELinux confines even uid-0 processes to the domains the
loaded policy permits, so that superuser identity gained within the
already-booted system does not by itself grant unrestricted operations
\citep{aosp-selinux}; the systemless root treated below escapes this
confinement not by wielding more privilege within the booted system but
by shaping the policy from beneath, as the boot image\textquotesingle s
init is patched before the policy is loaded.

The kernel\textquotesingle s lockdown facility, once engaged, withdraws
a range of otherwise root-accessible interfaces that could be used to
modify the running kernel \citep{kernel-lockdown}. Module-signature
enforcement can require that any code loaded into the kernel carry a
valid signature, so that root cannot insert an arbitrary kernel module
\citep{kernel-module-signing}, and the kernel.modules\_disabled control,
once set, blocks all further module loading until reboot in a one-way
fashion that even root cannot reverse \citep{kernel-modules-disabled}.

Beneath the operating system, Android Verified Boot and a locked
bootloader anchor the boot chain to a hardware root of trust and refuse
unsigned images, so that root -\/- an in-system privilege -\/- cannot by
itself rewrite the firmware or boot chain that sit below it
\citep{aosp-verified-boot}, \citep{aosp-bootloader-locking}. Throughout
this paper, therefore, "root" names a powerful but bounded in-system
privilege, not unconstrained control of the kernel, firmware, or boot
chain; Section 10 and the discussion in Section 14 return to why
obtaining it still does not deliver owner-inspectability.

One caveat attaches to the mechanisms just named:
Android\textquotesingle s Generic Kernel Image does not enforce a valid
signature on every loadable kernel module -\/- module signing is not
supported for GKI vendor modules, and on devices that require verified
boot the authenticity of kernel modules rests instead on dm-verity over
the read-only partitions that hold them
\citep{aosp-loadable-kernel-modules} rather than on per-module signature
enforcement. The mechanisms above are capabilities of the platform,
deployed to differing degrees across devices, not uniform guarantees on
every device.

Administrative control, by contrast, denotes the broader authority an
owner traditionally holds over a general-purpose computer: the ability
to replace the bootloader, kernel, kernel modules, initramfs, and
operating system, and to boot an external trusted environment of the
owner\textquotesingle s choosing, with uid-0 inside a system so
constituted. On a traditional personal computer this authority over the
operating-system and boot stack is the owner\textquotesingle s by
default; on stock Android it is withheld from the owner and conferred
only on the vendor. Root, when an owner obtains it at all, is a fragment
of this authority -\/- superuser standing that observes the system only
from within it, without command of the kernel and boot chain beneath it,
and in its prevalent systemless form obtained by patching init rather
than held as an account inside the booted system -\/- which is why, as
Sections 8, 10, and 14 develop, obtaining root does not amount to
administrative control and does not by itself deliver
owner-inspectability. Section 14 contrasts the two stacks directly,
including the firmware layers that remain opaque to the owner on both.

A clean, trusted environment denotes a separate computer or externally
booted system under the examiner\textquotesingle s control, as opposed
to the device\textquotesingle s own currently running operating system.

Owner-inspectability is the central property this paper isolates. A
device is owner-inspectable when its owner can examine and verify the
whole system -\/- storage, memory, and the boot and operating-system
state -\/- from a clean, trusted vantage point, and can establish and
demonstrate what the device is running, both to themselves and to any
verifier they choose to trust. The criterion is owner-facing: it does
not require the device to also keep the vendor-approved attested status
on which a remote relying party, not the owner, depends. That
vendor-approved standing is a property in tension with
owner-inspectability rather than a part of it -\/- on stock Android the
two are coupled, so that gaining owner inspection forfeits the
vendor-approved status, and it is that bind, not any inability to
self-verify, that this paper traces.

The term is deliberately distinct from forensic readiness as defined in
the digital-forensics literature, which concerns an
organization\textquotesingle s advance preparation to collect and
preserve admissible evidence about its own systems and users
\citep{rowlingson-forensic-readiness}, \citep{iso-27043}: forensic
readiness is a posture an institution arranges for investigators,
whereas owner-inspectability is a property the device either affords its
owner or withholds. Owner-inspectability is also narrower than a general
right to modify or repair a device: it requires that the owner be able
to establish, and to demonstrate, what the device is running -\/- not
that the device be freely modifiable. That separation is conceptual: in
practice the two entangle, since establishing what a locked device runs
-\/- most of all its live, whole-system state -\/- can itself require
modifying it (Section 14).

That separation is clean only for persistent, at-rest state: inspecting
the files, installed packages, and boot chain a device holds at rest is
a read-only act that requires no modification. Detecting non-persistent,
memory-resident malware is different, because such code leaves no
durable artifact and can be observed only through live memory
acquisition or runtime and kernel instrumentation
\citep{sudhakar-fileless-malware}, which on a locked device is itself a
modification -\/- unlocking, rooting, or booting a custom kernel
\citep{yang-live-memory-acquisition} -\/- that, as Section 8 shows,
forfeits the device\textquotesingle s attested status. For such live
threats inspectability and modifiability cannot be cleanly separated,
and the owner\textquotesingle s bind deepens rather than eases. AOSP in
principle allows a device to be locked to a user-controlled root of
trust rather than only to the vendor\textquotesingle s
\citep{aosp-verified-boot}; whether that latent capability is usable in
practice, and whether it survives contact with the attestation
ecosystem, is examined in Sections 8 and 12.

\subsection{4. Related Work}\label{related-work}

The platform mechanisms examined here are documented authoritatively in
the Android platform security model, which sets out verified boot,
application sandboxing, mandatory access control, and attestation as a
coherent design and states the threat model they address
\citep{mayrhofer-android-security-model}. That account is written from
the platform\textquotesingle s standpoint -\/- the device defending
itself, its users\textquotesingle{} data, and the ecosystem against a
running adversary. The present paper adopts the complementary standpoint
of the owner who wishes to audit the device, and reads the same
mechanisms for what they deny that owner.

Three bodies of prior work supply the evidence reorganized here. The
mobile-forensics literature establishes how, and how far, data can be
acquired from Android devices: official guidance on sound acquisition
and its preconditions \citep{nist-sp800-101r1}, empirical studies of
imaging and rooting-based acquisition
\citep{son-android-data-integrity}, \citep{vidas-android-collection},
and field reports on the limited traces recoverable from stock devices
in spyware investigations \citep{amnesty-pegasus-methodology},
\citep{mvt-android-methodology}. The rootkit-detection literature
establishes that trustworthy detection of kernel-level compromise
requires a vantage point outside the running operating system, and
explores hardware-assisted acquisition to obtain one
\citep{guri-joker-jtag}, \citep{bickford-rootkits-smartphones},
\citep{brodbeck-covert-rootkit-detection}. The anti-forensics and
evasion literature establishes that malicious code routinely detects
analysis environments and removes its own traces
\citep{chen-anti-virtualization}, \citep{ruggia-evasive-malware},
\citep{garfinkel-anti-forensics}. Each cluster is examined in detail in
the sections that follow; the contribution of this paper is not new
measurement within any one of them but their systematization around
owner-inspectability.

The owner-inspectability criterion relates to, but is distinct from,
organizational forensic readiness \citep{rowlingson-forensic-readiness},
\citep{iso-27043} and from remote-attestation design
\citep{mayrhofer-android-security-model},
\citep{play-integrity-improvements}; the distinctions are drawn where
each is used. To the author\textquotesingle s knowledge, the limits of
owner-side verification on stock Android have not previously been
assembled into a single, criterion-driven account.

Two further strands bear directly on the objections in Section 12. The
security cost of rooting as actually deployed has been measured: a study
of the commercial one-click-root ecosystem found that root providers
distribute kernel exploits that malware can equally abuse, so root as
shipped tends to enlarge the device\textquotesingle s attack surface
rather than give the owner controlled privileged access
\citep{zhang-root-providers}. And the tension between owner control and
remote attestation has a prior line of remedy: property-based
attestation proposes certifying that a platform satisfies required
security properties rather than that it runs one specific binary
configuration, which would let an owner-modified system attest without
being locked to a vendor image \citep{sadeghi-property-attestation}; the
owner-signed configuration examined in Section 12 is a concrete instance
of the same impulse, bounded in practice by attestation roots the owner
does not control.

\subsection{5. Independent Verification
Limits}\label{independent-verification-limits}

On most stock devices it is difficult to independently verify the state
of the whole system. When malware, a hidden backdoor, or tampering is
suspected, the owner generally cannot examine the system from a clean,
trusted environment without first modifying the device (Section 10). The
owner is therefore left to trust the vendor-controlled software and
services that deliver applications and updates, even though those
channels could themselves be abused in a targeted manner (Section 6).

The underlying cause is architectural. Android is commonly shipped as a
locked-down general-purpose computer in which administrative capability
is deliberately withheld from the owner and the bootloader is locked.
Many protections and integrity checks are implemented under vendor
control, limiting what the owner can change, inspect, or replace. As a
result the owner is typically prevented from inspecting the whole system
and from producing a complete device backup without first bypassing
these restrictions.

These restrictions bear on the three goals defined above (Section 3).
First, deep inspection for malware and backdoors is impeded: on a
non-rooted device, the private files of other applications are
inaccessible, and system areas and logs are access-controlled rather
than openly readable \citep{mayrhofer-android-security-model},
\citep{aosp-selinux}, which makes it harder to confirm whether
suspicious behavior originates from malware or from a normal system
component, and harder to rule out targeted tampering (Section 10).
Second, the owner cannot readily prove a clean or compromised result to
a third party, because the evidence that would support such a proof
cannot be collected. Third, a complete device backup -\/- for example, a
full offline image of the relevant partitions -\/- generally requires
privileges the device restricts (Section 7); on a non-rooted device
"full backup" in practice means a partial backup mediated by the
operating system and by individual applications
\citep{android-allowbackup}, \citep{aosp-auto-backup}.

For the same reasons, demonstrating the presence or absence of a hidden
backdoor is difficult. An examiner who cannot fully inspect system
files, logs, and low-level components without first modifying the device
is restricted to what the running operating system reports -\/- the same
limited vantage point a well-concealed backdoor would rely on (Section
9).

\subsection{6. Risks of Centralized Application and Update
Delivery}\label{risks-of-centralized-application-and-update-delivery}

Because key parts of application and update delivery are controlled by
the platform operator and device vendor -\/- on Android, Google Play for
applications and the manufacturer for system firmware -\/- vendor-locked
devices are exposed to the risk of targeted malicious application
installations, application upgrades, and system upgrades.

A documented and ordinary feature illustrates the mechanism. Many
devices permit a user to sign in to the Google Play web or desktop store
using the same Google account, typically the same email address, that is
used on the device. Selecting "Install" for an application in the web
interface prompts the user, where multiple devices are registered, to
choose a target device; the application is then installed on that
device. This demonstrates that Google Play can initiate an installation
to a specific, linked device without direct physical interaction with
that device \citep{google-play-help}.

The concern is not that an arbitrary outsider can invoke this capability
without authorization. The concern is that control of the delivery
system confers significant capability. If Google Play\textquotesingle s
infrastructure were compromised, or if a privileged insider were
malicious or coerced, this capability could in principle be abused to
push a backdoored version of an application to a specific target. By the
same logic, device manufacturers control firmware and system updates, so
targeted malicious system upgrades are likewise possible if that control
is abused.

Update-related communication is, moreover, not anonymous on devices
certified for Google Mobile Services (GMS). Google Play services
periodically checks in with Google, transmitting persistent hardware
identifiers -\/- the IMEI, hardware serial number, network MAC address,
and Google Services Framework identifier -\/- that Google documents as
associated with the signed-in Google account and used, among other
purposes, to determine a device\textquotesingle s eligibility for
software updates and security patches \citep{leith-handset-privacy},
\citep{google-device-config-service}. Whether a particular firmware
update is delivered through Google or through the device vendor varies
by manufacturer. The AOSP security policy itself grants Google Play
services access to the system-update service and to the A/B update
engine, confirming its privileged role in the update path on such
devices \citep{aosp-sepolicy-gmscore}. The data Google Play services
sends to Google in this way is not something the user can opt out of
while the service is in use: an independent measurement study of
GMS-certified handsets found that Google Play services and the Play
store transmit a considerable volume of data to Google even when the
device is idle and privacy settings are enabled, and reports that Google
confirmed there is no opt-out from this collection
\citep{liu-android-oem-snooping}.

Network-level anonymization does not neutralize this exposure, because
the exposure is at the identity layer rather than the network layer.
Using Google Play on a stock Android device requires the device to be
signed in to a Google account; the store cannot be browsed or used to
install applications without one \citep{google-play-help}. Every
interaction with Google Play authenticates that account to Google, so
each install or update request is bound to a specific, persistent
identity. Routing the device\textquotesingle s traffic through Tor
conceals the network (IP) address from which a request originates, and
in that narrow sense remains functional; but the account-based
authentication required by the operating system still presents the
user\textquotesingle s Google identity to Google on every request. The
network layer is thus anonymized while the identity layer is not. An
account identifier is not, by definition, the user\textquotesingle s
legal identity, but on a stock Android device it is in practice strongly
linkable to one: account creation commonly requires phone verification,
which Google documents as a check that the holder is a real person
\citep{google-verify-account}, and a phone number is itself bound to a
registered identity under the near-ubiquitous mandatory SIM-registration
regimes \citep{donovan-sim-registration}. The account also accumulates
an extensive, identity-linked record of activity and device identifiers
over time \citep{schmidt-google-data-collection}, and the network origin
from which it is used remains recoverable from provider-side logs, whose
retention many jurisdictions have mandated
\citep{eu-data-retention-directive}. Presenting the signed-in Google
account on each request therefore discloses an identifier strongly
linkable to the user, so routing the device\textquotesingle s traffic
through Tor conceals the network origin without conferring requester
anonymity.

\subsection{7. Limitations of Internal Storage Access on Non-Rooted
Devices}\label{limitations-of-internal-storage-access-on-non-rooted-devices}

Accessing internal storage for purposes such as data recovery or malware
and rootkit analysis presents significant practical barriers on most
stock devices. The guiding principle of malware analysis and backdoor
detection is that a device suspected of compromise should not be relied
upon to examine itself \citep{rfc3227}; the data should ideally be
inspected from a separate, known-clean computer, in keeping with
established forensic guidance on sound acquisition and evidence
preservation \citep{nist-sp800-101r1}. Stock Android prevents exactly
this.

The obstruction has two roots. First, internal storage is typically a
soldered component that cannot be removed without specialized work
\citep{nist-sp800-101r1}. Second, Android normally blocks access to many
system files and to the private files of other applications unless the
device is rooted or unlocked \citep{mayrhofer-android-security-model},
\citep{son-android-data-integrity}.

Reading the storage with a different device is generally impractical for
routine verification and is often unsuitable in forensic contexts. If
the device is physically damaged or unbootable, the storage is also
unlikely to function reliably when transplanted. Even where specialized
extraction is possible -\/- for example, chip-level desoldering and
reading -\/- it demands equipment and expertise that most owners do not
possess \citep{nist-sp800-101r1}.

The available software paths are correspondingly limited. The
platform\textquotesingle s backup facility covers only applications that
permit backup: it is enabled through the application manifest and
applications that disable it are omitted from the resulting archive, and
the mechanism is further bounded by per-application size limits and
excluded directories \citep{aosp-auto-backup}. Approaches that would
copy a full raw disk image to a computer do not work on devices that are
neither rooted nor rootable, and comprehensive imaging of internal
storage is generally described in the forensic literature as requiring
root or an unlocked bootloader \citep{son-android-data-integrity},
\citep{vidas-android-collection}.

These limitations are compounded by storage encryption, which is
effectively universal on recent mainstream devices: file-based
encryption has been mandatory on devices that launched with Android 10
or later, and has been available since Android 7
\citep{aosp-file-based-encryption}. The consequences of a hardware-bound
encryption key are examined further in Section 11. Encryption is a
second barrier behind the first, not the main one. The main obstacle to
offline imaging is physical: for most owners the soldered chip cannot be
removed and read at all, as described above. Where chip-level extraction
does nonetheless succeed, encryption limits what it yields -\/- the
user-data partition comes out as ciphertext, because the key that
decrypts it is derived on-device from a secret bound to hardware-backed
secure storage and cannot be reproduced off the device by routine
extraction. The read-only system partitions are not encrypted this way
and can be read, but they hold the vendor\textquotesingle s code rather
than the owner\textquotesingle s data. The owner\textquotesingle s own
data is therefore doubly out of reach: the chip usually cannot be read,
and on the rare occasion it can, that data is ciphertext. For most
owners, therefore, inspection of internal storage is not reasonably
achievable on a stock, non-rooted device.

\subsection{8. Attestation and the Circular Trust
Problem}\label{attestation-and-the-circular-trust-problem}

Stock Android relies heavily on remote device attestation -\/-
historically Google\textquotesingle s SafetyNet Attestation API, which
Google deprecated in 2022 and withdrew in January 2025
\citep{safetynet-attestation-deprecation}, now replaced by the Play
Integrity API \citep{play-integrity-overview} -\/- by which the device
cryptographically reports aspects of its own state to applications and
services. The reported state is grounded in the device\textquotesingle s
verified-boot status \citep{aosp-verified-boot} and backed by
hardware-backed key attestation \citep{aosp-key-id-attestation}. Table 1
summarizes three simplified device-state claims that the platform
exposes through its device-integrity signals; the Play Integrity API
additionally returns separate application-integrity and
account-integrity fields that are not reproduced here
\citep{play-integrity-overview}.

Table 1. Device attestation: what is reported to applications.

\begin{longtable}[]{@{}
  >{\raggedright\arraybackslash}p{(\columnwidth - 2\tabcolsep) * \real{0.5000}}
  >{\raggedright\arraybackslash}p{(\columnwidth - 2\tabcolsep) * \real{0.5000}}@{}}
\toprule\noalign{}
\begin{minipage}[b]{\linewidth}\raggedright
Attestation claim
\end{minipage} & \begin{minipage}[b]{\linewidth}\raggedright
Meaning (simplified)
\end{minipage} \\
\midrule\noalign{}
\endhead
\bottomrule\noalign{}
\endlastfoot
Running vendor-approved software & The device reports that it runs an
operating system or build approved by the vendor, and often by
Google. \\
Unmodified & The device reports that it has not been altered, for
example that the bootloader has not been unlocked and the system has not
been modified. \\
Locked down & The device reports that security-relevant settings remain
enforced, including a locked boot chain and an expected integrity
state. \\
\end{longtable}

Attestation is a relying-party protocol, and naming its roles makes the
direction of trust explicit. In the standard architecture the device is
the Attester, which must produce evidence about itself; a Verifier
appraises that evidence; and a Relying Party -\/- here the application
developer, content provider, or platform operator -\/- is, in that
architecture\textquotesingle s words, the party that "depends on the
validity of information about an Attester" before it will grant service
\citep{rfc9334-rats}. Remote attestation is in general the activity of
making a claim about a target\textquotesingle s properties by supplying
evidence to an appraiser over a network
\citep{coker-remote-attestation}.

The guarantee therefore runs to the relying party, not to the device
owner: the device proves its compliance to a remote party that does not
extend trust to the owner, and a passing verdict certifies the device to
that party rather than giving the owner any means to audit or halt what
an application does. That inversion -\/- software operating on the
owner\textquotesingle s own device while it is the owner who is kept
from inspecting or stopping it -\/- has been characterized as a broader
loss of owner control over general-purpose computing
\citep{doctorow-general-computation}. Accordingly, both "trusted status"
and "attested status" in this paper denote this attested,
vendor-approved standing that a remote party relies on, not any trust
the owner places in the device.

Android does surface a status to the owner -\/- the Play
Store\textquotesingle s Play Protect certification
\citep{play-protect-certification} and the verified-boot warnings shown
at startup \citep{aosp-boot-flow} -\/- but that status reports the
platform\textquotesingle s own approval verdict, reflecting a locked
bootloader and a vendor-certified image; the owner must trust that
verdict rather than independently audit it, and it is not a verification
that the device is actually uncompromised: the device-integrity verdict
attests boot and configuration state, not runtime malware-freeness
\citep{play-integrity-verdicts}, so a locked, certified device can still
be compromised in ways this boot-time check does not cover -\/-
malicious code shipped inside the certified image as pre-installed
vendor software, as the Triada family was built into
manufacturers\textquotesingle{} system images during production
\citep{google-triada}, or a runtime, memory-resident compromise of the
already-booted and legitimately signed kernel, such as the in-the-wild
Binder use-after-free in the Android kernel that escalated to full
device compromise \citep{projectzero-badbinder}.

These runtime compromises are documented, not hypothetical: commercial
surveillance vendors sell turnkey exploit chains that defeat current,
fully patched devices, and Google\textquotesingle s Threat Analysis
Group attributes about half of known zero-day exploits against Google
and Android targets to such vendors \citep{tag-buying-spying}, in one
case observing a full chain fired against an up-to-date Samsung phone
running the latest Chrome \citep{tag-cytrox}; Project Zero maintains the
running in-the-wild record \citep{projectzero-itw}. Once resident, such
an implant operates inside the running system\textquotesingle s own
trust boundaries -\/- the Predator spyware\textquotesingle s loader
injects into zygote to assume other application identities and move
between SELinux domains \citep{talos-predator} -\/- so the application
sandbox and SELinux confinement of Section 2 restrain it far less than
they restrain ordinary software.

A persistent rootkit that instead rewrites the verified boot or system
partitions is the one case verified boot and dm-verity are built to
catch \citep{aosp-verified-boot}; Section 12 returns to the deeper
breaches that can defeat even that.

This arrangement produces a circular trust dependency, summarized in
Table 2. To inspect the system deeply, the owner must modify the device
-\/- unlock the bootloader, obtain root, or install special builds -\/-
which causes the device to fail attestation and to be treated as
untrusted by applications. To remain trusted by applications, the owner
must keep the device unmodified and locked down, which denies access to
the disk, memory, logs, and low-level components that independent
verification requires.

Table 2. Circular trust dependency: inspection versus attestation.

\begin{longtable}[]{@{}
  >{\raggedright\arraybackslash}p{(\columnwidth - 4\tabcolsep) * \real{0.3333}}
  >{\raggedright\arraybackslash}p{(\columnwidth - 4\tabcolsep) * \real{0.3333}}
  >{\raggedright\arraybackslash}p{(\columnwidth - 4\tabcolsep) * \real{0.3333}}@{}}
\toprule\noalign{}
\begin{minipage}[b]{\linewidth}\raggedright
Goal
\end{minipage} & \begin{minipage}[b]{\linewidth}\raggedright
What is typically required
\end{minipage} & \begin{minipage}[b]{\linewidth}\raggedright
Consequence
\end{minipage} \\
\midrule\noalign{}
\endhead
\bottomrule\noalign{}
\endlastfoot
Inspect the system deeply & Modify the device (unlock bootloader, gain
root, use special builds) & Device may fail attestation and be treated
as untrusted by applications \\
Remain trusted by applications & Keep the device unmodified and locked
down & Limits access to the disk, memory, logs, and low-level components
needed for independent verification \\
\end{longtable}

From the platform\textquotesingle s standpoint this is a security
feature: it raises assurance that a device presenting itself to a
service is in a known, unmodified state \citep{play-integrity-verdicts}.
From the standpoint of an owner facing a high-threat situation, the same
mechanism forecloses independent verification, because the only
configurations in which the device is considered trustworthy by the
applications and services that demand it are precisely those in which
the owner cannot examine it.

The dependency has, moreover, tightened over time, and now reaches
beyond the owner\textquotesingle s own verification into everyday use of
the device. Attestation verdicts are tiered: the platform distinguishes
a basic verdict from a device-integrity verdict and a stronger,
hardware-backed verdict, and it exposes which tier a given device can
satisfy \citep{play-integrity-improvements}. A device whose owner has
unlocked the bootloader, installed an owner-controlled operating system,
or obtained root typically cannot satisfy the device-integrity or strong
verdicts, because those are tied to a Google-certified, vendor-signed
configuration \citep{play-integrity-improvements},
\citep{owasp-mastg-play-integrity}. Some applications -\/- banking and
payment applications and content-protection components among them -\/-
gate access on at least the hardware-backed device-integrity verdict,
some additionally demanding the strong tier, and refuse to run, or
silently degrade, on a device that cannot present them
\citep{owasp-mastg-play-integrity}, \citep{play-integrity-verdicts},
\citep{aosp-drm}, \citep{privsecdev-banking-grapheneos}.

The exact proportion of applications that do so is disputed and not
rigorously measured. The available evidence is limited and indirect, but
the measurements that exist are consistent with refusal being common and
increasing. A code-level measurement reports that most banking
applications contain no reference to Google\textquotesingle s own
attestation frameworks in their first-party code -\/- roughly nine in
ten reference neither the Play Integrity API nor the deprecated
SafetyNet, and the minority that do still overwhelmingly ship only
SafetyNet calls \citep{macdonald-attestlens}. That figure is a floor on
lock-out, not a measure of it: such a count detects only
Google\textquotesingle s own frameworks, and the same measurement
observes that most attestation that is present originates from
third-party packages rather than first-party code, so intentional
adoption is systematically undercounted \citep{macdonald-attestlens}.

The enforcement an owner actually encounters is delegated elsewhere -\/-
to commercial application-shielding and runtime application
self-protection (RASP) SDKs such as Promon SHIELD, Guardsquare/DexGuard,
DexProtector, and Arxan, and to bespoke root, hook, and emulator checks
that invoke no attestation API at all. A peer-reviewed systematization
of mobile hardening finds these pervasive on current Android: root
detection in 92.3 percent of apps, debugger detection in roughly 94
percent, and emulator detection in nearly 99 percent, with hardening
most concentrated in the finance and shopping categories and, by that
study\textquotesingle s own account, driven largely by third-party
libraries \citep{steinboeck-sok-hardening}. Consistent with this,
crowd-sourced compatibility lists for owner-controlled operating systems
catalog a set of banking and payment applications that refuse to run or
install on them \citep{privsecdev-banking-grapheneos},
\citep{grapheneos-attestation-compat}.

The summary is therefore a floor-and-direction one: the first-party Play
Integrity figure is a lower bound, the delegation of enforcement to RASP
SDKs and non-attestation checks is why real-world refusal is
substantially more common than that bound, and no single measurement yet
fixes a precise lock-out rate. The owner, meanwhile, has no way to
appeal a verdict except by returning the device to the vendor-approved
state. The consequence for this paper\textquotesingle s argument is
structural: even where an owner is technically able to reach an
owner-inspectable configuration, the attestation-and-hardening ecosystem
imposes a rising functional cost for remaining there, so that the
circular trust dependency is enforced not only by the
platform\textquotesingle s own refusal to trust a modified device but by
the third-party software the owner must run on it.

This also sharpens what "gaining root" can and cannot achieve, a point
developed in Sections 10 and 14. Because root is an in-system privilege
(Section 3) and because reaching it generally forfeits the attestation
standing just described, obtaining root does not hand the owner the
clean, external, trusted vantage point that independent verification
requires; it trades the device\textquotesingle s attested status for an
in-system foothold that still observes the kernel only from within it
-\/- and, where lockdown or module-signature enforcement is engaged,
cannot even load code into that kernel.

Modern rooting fits this bound rather than breaking it. The prevalent
method is systemless: it gains root without writing to /system. Verified
boot covers /system (guarded by a dm-verity hashtree
\citep{aosp-dm-verity}) and the boot and init\_boot images (an AVB hash
in the signed vbmeta) alike, so on a locked bootloader neither can be
altered; the unlock is the shared precondition, after which the device
boots an unsigned image and Magisk disables vbmeta verification
\citep{aosp-verified-boot}, \citep{magisk-install}. It patches the init
carried in the ramdisk -\/- in the boot partition, the init\_boot
partition on devices launched with Android 13 or later, or, on some
devices, vendor\_boot or a recovery image \citep{aosp-partitions},
\citep{magisk-install} -\/- because init is the earliest userspace hook,
then overlays its changes at runtime, leaving the read-only /system
byte-identical: official OTAs stay applicable (root is re-applied to the
updated slot afterward, not carried across) and the change is reversible
by reflashing the stock boot image \citep{magisk-install},
\citep{magisk-ota}. Modern "root" is in this sense better described as
control of init than as an in-system account: the patched init runs
before the SELinux policy is loaded, so the installer shapes the policy
its own daemon then runs under -\/- an unconstrained, permissive domain
confined to that daemon and the root shells it grants, while every other
domain stays enforcing \citep{magisk-details} -\/- rather than being
confined by it.

Even so, this foothold is not the external vantage point verification
needs. It is the administrative-tier power to replace the initramfs
(Section 3), reached only through the unlock that forfeits verified
boot, and the permissive domain sidesteps SELinux but not the running
kernel beneath it: loading code into that kernel is blocked where
lockdown, module-signature enforcement, or the one-way module-disable
control are engaged \citep{kernel-lockdown},
\citep{kernel-module-signing}, \citep{kernel-modules-disabled}, though
not uniformly -\/- the Generic Kernel Image does not enforce per-module
signatures (Section 3) -\/- and in any case this root observes the
kernel only from within it, so a kernel-space rootkit beneath it
falsifies what any in-system tool reports regardless of privilege
(Sections 9 and 14).

A clean, out-of-system vantage comes not from the ramdisk patch but from
the same unlock\textquotesingle s further use -\/- booting a separate
kernel or recovery to inspect the partitions at rest -\/- and even that
leaves the suspect kernel\textquotesingle s live memory beyond trusted
reach (Section 14). The modification has thus moved down into the boot
chain that verified boot protects, not out of its reach, and the device
still forfeits the hardware-backed verdicts described above.

The integrity-masking modules an owner adds alongside it spoof the
software-visible signals read by app-side root checks and the basic
verdict, and on devices launched before Android 13 -\/- whose
device-integrity verdict is the weaker successor to
SafetyNet\textquotesingle s ctsProfileMatch -\/- they can pass that
verdict too by rewriting android.os.Build fields and system properties,
with no hardware key material \citep{playintegrityfork}. What software
spoofing cannot forge is the hardware-backed device-integrity verdict
that devices launched with Android 13 and later enforce through
locked-bootloader checks, nor the strong verdict: passing those requires
genuine certified hardware with a locked bootloader, or leaked,
not-yet-revoked hardware key material of the kind Section 12 discusses
\citep{play-integrity-verdicts}, \citep{play-integrity-improvements},
\citep{grapheneos-attestation-compat}, \citep{playintegrityfork} -\/- a
key compromise rather than a defeat of the mechanism.

\subsection{9. Detection of Low-Level
Rootkits}\label{detection-of-low-level-rootkits}

Low-level rootkits conceal themselves by manipulating the operating
system itself, so that security tools running on the same device may be
unable to detect them reliably. Trusted detection can require obtaining
a trusted snapshot of memory from outside the monitored operating system
\citep{guri-joker-jtag}.

The problem is structural rather than a matter of tool quality.
Anti-virus software may successfully detect malicious applications yet
remain ineffective against low-level rootkits that evade detection by
masking their own presence \citep{guri-joker-jtag}. Any detection
mechanism that runs on the same physical device as the monitored
operating system can itself be compromised through application, kernel,
or bootloader vulnerabilities, so trusted detection of kernel rootkits
on mobile devices is, in practice, a challenging task
\citep{guri-joker-jtag}.

The distinction between user-space and kernel-space malware clarifies
why. User-space malware can modify or inject code only into memory
allocated to applications and user processes. Kernel-space malware can
manipulate objects throughout the operating system\textquotesingle s
memory. Although mandatory access-control mechanisms such as SELinux are
integrated into current Android versions, rootkits running at the
operating-system level of mobile devices have been demonstrated
\citep{bickford-rootkits-smartphones},
\citep{brodbeck-covert-rootkit-detection}. A rootkit is kernel-space
malware that uses illicitly obtained privileges to hide its existence by
manipulating the kernel\textquotesingle s internal data structures, and
code that has penetrated kernel memory can neutralize any security tool
running within the operating system. A request to the kernel for, say,
the contents of a directory carries no guarantee of integrity in the
returned list; consequently, detecting a rootkit requires a trusted
snapshot of kernel memory obtained from outside \citep{guri-joker-jtag}.

One proposed approach obtains such a snapshot through the JTAG hardware
debugging interface rather than through the potentially compromised
operating system. The approach exploits two debugging features: the
ability to halt the processor instantly by issuing a special
instruction, and the ability to read the device\textquotesingle s
volatile memory while it is halted. Because the acquisition system does
not execute on the device, it can read kernel memory in a trusted
manner; the captured memory is then processed by scripts that
reconstruct kernel data structures and analyze them for suspicious
modifications. The acquisition offers the advantage of trusted
inspection precisely because it originates from outside the monitored
system \citep{guri-joker-jtag}.

For most owners and most devices, JTAG-based memory acquisition is not
practical. Some devices expose a JTAG interface, but it is usually not
accessible to typical owners; access tends to be device-specific and to
require specialized equipment and skills. No quantitative data on the
proportion of devices for which JTAG is usable has been located, and the
situation is highly device-specific. Manufacturers may treat JTAG as an
attack surface and disable or lock it on production hardware
\citep{patent-secure-jtag}, \citep{sguigna-jtag-attack-surface}: a
recommended practice is to disable JTAG by default in production and
end-user devices to reduce attack surface, enabling it only temporarily
for authorized maintenance \citep{sguigna-jtag-attack-surface}. The same
interface can nonetheless serve security analysis, and is an established
forensic acquisition method for reading memory from embedded systems
\citep{breeuwsma-jtag-imaging}. The approach\textquotesingle s own
authors acknowledge its practical limitations: it requires physical
access to a JTAG port on the circuit board, which is cumbersome compared
with software-based methods \citep{guri-joker-jtag},
\citep{brodbeck-covert-rootkit-detection}.

On many devices, deep investigation -\/- including the search for
sophisticated rootkits -\/- requires root access, and where the
bootloader cannot be unlocked, researchers may be driven to rely on
exploitable vulnerabilities to gain privileges temporarily. Forensic
analyses report that application sandboxing, encryption, and file-based
access controls pose challenges to data acquisition
\citep{nist-sp800-101r1}, and that rooting -\/- commonly achieved by
executing an exploit against an operating-system vulnerability, or by
booting a custom recovery mode that grants administrative privileges
while minimizing alteration of stored user data -\/- provides the access
needed for fuller acquisition \citep{son-android-data-integrity},
\citep{vidas-android-collection}.

For comparison, non-Android platforms such as desktop and server Linux
provide options for the same objectives. Anti-malware kernel modules
designed to detect loadable-kernel-module rootkits, and runtime
integrity guards, can run on the system \citep{thalium-rkchk},
\citep{openwall-lkrg}, though in-system monitors remain subject to the
same kernel-level evasion described in Section 9. Alternatively, the
system can be booted from external media and subjected to an offline
integrity check using tools that verify installed files against a
trusted manifest \citep{elstel-debcheckroot}. These options permit
independent verification even where the original operating-system
installation is itself suspected.

\subsection{10. Malware-Analysis Capability
Gaps}\label{malware-analysis-capability-gaps}

The constraints above translate into concrete gaps in analysis
capability. This section contrasts a fully controlled general-purpose
host, represented by traditional Linux desktops and servers, with a
stock Android device.

On a traditional Linux host the examiner can replace the operating
system, installing a known-good system or reinstalling from trusted
media if compromise is suspected; boot from external media to obtain a
clean, trusted environment without relying on the suspect system disk;
inspect disks and memory from that trusted environment while minimizing
reliance on the suspect operating system; instrument or replace the
kernel to load debugging or integrity tooling; and run security tools
outside the potentially compromised system, from an external host or
trusted boot environment. Taken together, these capabilities permit
independent verification even when the original installation is
suspected to be malicious.

On a stock Android device these same capabilities are, in general,
absent. The operating system cannot be freely replaced, so restoring a
known-good system is usually unsupported or impractical. Booting an
external, fully trusted operating system that bypasses the installed
system is not supported on the device as sold: while the bootloader is
locked it requires a valid signature before it will boot or install any
image, so an external, unsigned system cannot be started without first
unlocking it \citep{aosp-fastboot-protocol},
\citep{aosp-bootloader-locking}. Even on an unlocked device there is no
universal external-boot path -\/- starting an external image is vendor-
and device-specific, and cumbersome even where it works at all
\citep{bootloader-wall-of-shame}. Full disk and memory access -\/- raw
disk imaging, comprehensive file inspection, and live RAM capture -\/-
is blocked without special privileges. Gaining those privileges usually
breaks verified boot and attestation, causing the device to be treated
as untrusted; can disable or degrade security features that depend on
the device remaining locked down and unmodified; and can cause many
applications, including banking, enterprise, and content-protected
applications, to refuse to run. The net effect is that the tools needed
to verify the system require breaking those same guarantees before they
can be applied -\/- and, as Section 8 showed, the privileges so gained
are in-system, not the external, trusted vantage point that verification
needs, and no amount of in-system privilege supplies one.

Table 3 compares specific analysis capabilities. The values "Yes" and
"No" indicate that a capability is, respectively, generally available or
generally unavailable to the examiner on a stock device; "Partial"
indicates that the capability exists only under significant constraints,
typically requiring root, debugging permissions, or vendor or forensic
access. The entries describe the device in its stock, as-sold
configuration. An owner who unlocks the bootloader and installs an
owner-controlled system can move several Android entries toward the
Linux column, but only by forfeiting the attestation standing discussed
in Section 8 and, with it, the ability to run the applications that
depend on it; the comparison therefore characterizes the trade-off the
owner faces on the device as sold, not an immutable property of the
hardware.

Table 3. Comparison of malware-analysis capabilities: a traditional
Linux host the owner administers versus stock Android (as sold).

\begin{longtable}[]{@{}
  >{\raggedright\arraybackslash}p{(\columnwidth - 4\tabcolsep) * \real{0.3333}}
  >{\raggedright\arraybackslash}p{(\columnwidth - 4\tabcolsep) * \real{0.3333}}
  >{\raggedright\arraybackslash}p{(\columnwidth - 4\tabcolsep) * \real{0.3333}}@{}}
\toprule\noalign{}
\begin{minipage}[b]{\linewidth}\raggedright
Aspect
\end{minipage} & \begin{minipage}[b]{\linewidth}\raggedright
Traditional Linux
\end{minipage} & \begin{minipage}[b]{\linewidth}\raggedright
Stock Android (as sold)
\end{minipage} \\
\midrule\noalign{}
\endhead
\bottomrule\noalign{}
\endlastfoot
System access (examiner-controlled) & Yes -\/- full root access common
in testbeds and forensic imaging & No -\/- root typically restricted;
requires manufacturer unlock (often unavailable) or an exploit to gain
privileged access \\
User-controlled root & Yes -\/- commonly available in lab and test
environments and on self-managed systems & No -\/- not available on the
device as sold; typically requires bootloader unlock, custom builds, or
an exploit \\
Raw disk imaging & Yes -\/- straightforward via standard forensic
imaging workflows, including booting from external media & No -\/- not
feasible on the device as sold; generally requires root, an unlocked
bootloader, or specialized vendor or forensic access \\
Live memory acquisition & Yes -\/- readily possible with standard
memory-capture tooling & No -\/- not feasible without root; no official
RAM-acquisition interfaces \\
Kernel analysis & Yes -\/- kernel modules can be inserted, traced, or
extracted; symbols often available & No -\/- unavailable on the device
as sold without root, special builds, or vendor support \\
Kernel integrity verification & Yes -\/- possible using external boot
media and host-controlled integrity tooling & No -\/- not available in a
user-controlled way; verification is typically limited and
vendor-controlled \\
Independent integrity verification & Yes -\/- possible by booting a
trusted environment and performing independent integrity checks & No
-\/- limited on the device as sold; attestation and integrity guarantees
are typically vendor-controlled \\
Binary reverse engineering & Yes -\/- straightforward access to malware
files, system binaries, and logs & Partial -\/- possible, but often
limited to application-level packages unless rooted or firmware access
is available \\
Hook and injection analysis & Yes -\/- library preloading, process
tracing, and syscall interposition easily observed & Partial -\/-
requires rooting; SELinux and sandboxing block injection methods by
default \\
Dynamic instrumentation & Yes -\/- tracing and debugging tools available
and effective & Partial -\/- requires root or debugging permissions \\
Filesystem analysis & Yes -\/- full access to logs, temporary files,
droppers, and rootkits & Partial -\/- the application sandbox limits
access to other applications\textquotesingle{} data; root or a forensic
image is required for a full view \\
Network monitoring and command-and-control analysis & Yes -\/-
packet-capture and filtering tools usable without restriction & Partial
-\/- requires root or VPN-based monitoring \\
External monitoring & Yes -\/- broad options, including out-of-band
approaches and host-controlled capture and analysis & No -\/- limited on
the device as sold; out-of-band monitoring is generally unavailable
without elevated access or specialized tooling \\
Log-based forensics & Yes & Partial -\/- limited; absence of traces
after thorough manual analysis does not establish that a device was not
compromised, and log-based methods may need to be combined with other
forensic techniques where compromise is strongly suspected
\citep{socialtic-log-forensics} \\
\end{longtable}

\subsection{11. Limitations on Encryption Key
Backups}\label{limitations-on-encryption-key-backups}

Backing up the device\textquotesingle s encryption keys presents
significant practical barriers. Recent devices use file-based
encryption, mandatory since Android 10
\citep{aosp-file-based-encryption}; the keys that protect user data are
not held as ordinary files the owner can copy but are wrapped by a key
bound to hardware-backed secure storage, such as a secure element or
trusted execution environment, which places the protective key beyond
the owner\textquotesingle s reach. (Earlier devices used full-disk
encryption, which likewise bound its master key to hardware rather than
exposing it as a copyable file \citep{aosp-full-disk-encryption}.) The
use of a hardware-protected key is not in itself anomalous or an
indication of a backdoor.

The underlying arrangement -\/- a master key that decrypts the volume,
itself unlocked by the user\textquotesingle s passphrase -\/- is the
normal one in Linux full-disk encryption. There the master key is stored
on disk wrapped by the passphrase, so entering the passphrase at boot
decrypts the master key, which in turn decrypts the volume. The
arrangement allows the passphrase to be changed without re-encrypting
the entire volume, since only the stored master key must be re-wrapped
\citep{luks-cryptsetup-faq}. The difference on a stock device is that
hardware protection of the key places it beyond the
owner\textquotesingle s reach rather than leaving it in a file the owner
could back up.

Backing up the key may be possible on some devices if the device still
boots and root can be obtained, but for most owners this is not
straightforward.

\subsection{12. Objections and
Counterarguments}\label{objections-and-counterarguments}

A principal objection holds that practical forensic tooling for mobile
devices already exists and therefore contradicts the claim that stock
devices resist verification. The Mobile Verification Toolkit (MVT), a
consensual forensic tool intended to help analyze mobile devices for
signs of spyware, is the example considered here.

There is, however, no contradiction. Such a tool can identify known
indicators of compromise using the data that the device exposes, but the
stock mobile operating-system design limits that observability. The
toolkit\textquotesingle s own documentation states that Android devices
provide much less observability than comparable iOS devices, that
Android stores very little diagnostic information useful for triaging
potential compromises, and that the toolkit\textquotesingle s Android
capabilities are correspondingly limited, while noting that several
forensic artifacts remain accessible \citep{mvt-android-methodology}.
Investigators are thus obliged to work with partial data and with known
indicators of compromise \citep{mvt-android-methodology}, and a lack of
findings is not definitive proof that a device was not compromised.

Independent forensic experience reinforces this. Investigating NSO
Group\textquotesingle s Pegasus spyware, Amnesty
International\textquotesingle s Security Lab reported that significantly
more forensic traces are accessible to investigators on Apple iOS
devices than on stock Android devices, and that its methodology was
focused on iOS for that reason; as of that 2021 report, most of the
then-recently confirmed Pegasus infections had involved iPhones
\citep{amnesty-pegasus-methodology}. The report cautions that this does
not by itself establish that iOS is more secure than Android
\citep{amnesty-pegasus-methodology}. The same report characterized the
difficulty of detecting attacks even after the fact as the result of an
asymmetry between the capabilities available to attackers and the
protections available to targeted individuals, observed that devices
other than iOS provide little to no help for consensual forensic
analysis, and encouraged device vendors to make their devices more
auditable without sacrificing existing security and privacy protections
\citep{amnesty-pegasus-methodology}. These observations are consistent
with limited observability being a property of the platform design.

A second objection points to aftermarket Android-based operating systems
-\/- GrapheneOS, CalyxOS, LineageOS, /e/OS, postmarketOS, and others
-\/- as counterexamples that an owner can install in place of the
vendor\textquotesingle s build. They do not form a single case, and the
differences are exactly what the owner-inspectability criterion
measures.

Most such systems are installed with the bootloader left unlocked,
because the device cannot be re-locked against a key the vendor does not
recognize, and an unlocked bootloader cuts both ways against the
owner-inspectability criterion. It restores much of what the
criterion\textquotesingle s first half demands: the owner regains
privileged, at-rest inspection a stock device withholds -\/- imaging the
partitions above all, and, on well-supported hardware, booting an
owner-controlled kernel or recovery to instrument the system -\/- the
capabilities Sections 9 and 10 found absent. This reaches less far in
practice than in principle: a fully functional owner-built kernel is
confined to a minority of devices, because the vendor
hardware-abstraction interface \citep{aosp-hal}, the generic-kernel
module boundary \citep{aosp-gki}, \citep{aosp-loadable-kernel-modules},
and non-rebuildable firmware blobs tie working hardware to the
vendor\textquotesingle s own code; and the suspect
system\textquotesingle s live memory stays out of reach, since live
acquisition runs in the already-booted kernel, so booting
one\textquotesingle s own kernel first reboots the device and discards
the volatile state that soldered memory already places beyond trusted
external capture (Section 14).

What it gives up is the criterion\textquotesingle s second half: without
a locked chain the boot images are no longer cryptographically verified
against a hardware root of trust on each startup, and although hardware
key attestation still functions it can now report only an unverified
boot state, not certify a locked, vendor-approved one, so a tamper of
the boot chain between inspections is no longer prevented but at best
detectable by hand. Such an installation therefore satisfies inspection
while giving up the enforced, attestable verified-boot state -\/- an
illustration of the bind this paper describes, which the
owner-inspectability criterion is designed to capture.

A smaller number of them support re-locking the bootloader after
installation \citep{grapheneos-install-locking},
\citep{grapheneos-setup-wizard}, and this is where the
owner-inspectability criterion discriminates most finely, because the
re-locked chain can be anchored to either of two very different keys.
Anchored to the project\textquotesingle s own key -\/- the default when
an owner flashes official GrapheneOS release builds, which replace the
vendor\textquotesingle s verified-boot key in the
device\textquotesingle s secure element with
GrapheneOS\textquotesingle s own \citep{aosp-device-state},
\citep{grapheneos-install-locking} -\/- re-locking restores verified
boot but only substitutes one vendor\textquotesingle s root of trust for
another\textquotesingle s: the signing key is now
GrapheneOS\textquotesingle s rather than Google\textquotesingle s, and
the owner still does not hold it. On the key-control half of the
criterion this official-build configuration is the stock position under
a different vendor\textquotesingle s key, not owner control of the boot
chain; replacing an official Google image with an official GrapheneOS
image changes the vendor, not the owner\textquotesingle s standing. Even
if Google allowlisted such a build so that it passed the hardware-backed
integrity verdicts that banking and payment apps gate on, that would be
the attestation ecosystem and its relying parties choosing to accept
GrapheneOS as trusted, not the owner inspecting or verifying the device.

GrapheneOS is one case on which the thesis can be tested, because it
adds what the stock platform withholds from the owner: owner-accessible
attestation tooling -\/- the Auditor app, which verifies, from a second
device or on a schedule against an attestation server the owner can run
themselves \citep{grapheneos-auditor},
\citep{grapheneos-attestation-server}, that the system is the expected,
untampered build running on a locked bootloader, built on the standard
hardware attestation primitive but packaged as a verify-to-yourself
capability rather than a relying-party one.

The configuration that actually meets the key-control half is the other
anchoring: an owner who builds GrapheneOS from source, signs it with
keys they generate and keep, and re-locks to their own enrolled
verified-boot key holds a root of trust they control
\citep{grapheneos-build}; tooling exists to produce such owner-signed
builds from signed over-the-air updates without a full source rebuild,
though only on supported Pixel hardware, and it too documents that no
owner-signed configuration can satisfy the hardware-backed Play
Integrity verdict \citep{resign-android-image}.

GrapheneOS\textquotesingle s own Auditor app, in its released form,
nonetheless does not recognize that owner key: its attestation accepts
only official GrapheneOS keys in the self-signed verified-boot state
-\/- an alternate OS re-locked under a non-stock key -\/- because an
alternate system\textquotesingle s verified-boot key must be compiled
into the released app and server to be recognized
\citep{grapheneos-auditor}, \citep{grapheneos-auditor-source},
\citep{grapheneos-auditor-usecase} (Appendix A), so verifying an
owner-signed build needs a verifier the owner configures with their own
key -\/- an owner-built Auditor and, if remote checking is wanted, an
owner-run attestation server. That owner-signed path is the
configuration that actually satisfies the owner-inspectability criterion
for the main system -\/- the owner both holds the root of trust and can
verify the device against it, to themselves or to any verifier they
provision -\/- and is the configuration this paper uses to test the
thesis.

Even there one limitation is relocated rather than dissolved, and it is
worth naming precisely. The owner-signed configuration delivers verified
boot under the owner\textquotesingle s own key and an attestation the
owner can check -\/- to themselves, or to any verifier willing to adopt
that key as a trust anchor -\/- so for the owner\textquotesingle s own
purposes, establishing and demonstrating what the device runs, both
halves of the criterion are met. What it cannot do is prove itself to a
verifier that refuses in advance to trust the owner\textquotesingle s
key: hardware attestation still demonstrates genuine hardware in a
locked, self-signed boot state to anyone who checks, but a skeptic has
no published, independent anchor for the owner\textquotesingle s
verified-boot key, so the gap is the absence of a trust anchor, not the
unavailability of the key -\/- publishing the key would not close it. An
official GrapheneOS key differs only in that a relying party can choose
to trust GrapheneOS, a public project with a reproducible build, as that
anchor; a self-generated key vouches for no one but its owner. For
everyday use, neither the owner-signed build nor the official GrapheneOS
build is a Google-certified configuration, so neither obtains the
device-integrity or strong attestation verdicts that mainstream
applications increasingly demand \citep{play-integrity-improvements},
\citep{owasp-mastg-play-integrity}: the same owner-controlled posture
that makes the device inspectable makes an apparently growing set of
applications refuse to run on it, and GrapheneOS\textquotesingle s
stated position is that it will not spoof or otherwise circumvent these
checks to impersonate a stock device \citep{grapheneos-usage}.

Other owner-controlled systems take the opposite tack: the microG-based
distributions CalyxOS and /e/OS bundle a reimplementation of Google Play
Services that depends on signature spoofing to impersonate those
services to the apps that call them \citep{microg-gmscore},
\citep{calyxos-microg}, \citep{eos-microg}, and rooted owners add
integrity-masking modules -\/- Universal SafetyNet Fix and the Play
Integrity Fix family, optionally paired with a keybox provider such as
Tricky Store -\/- that spoof the attestation signals themselves
\citep{safetynet-fix}, \citep{playintegrityfork}, \citep{tricky-store}.
Signature spoofing of this kind restores access to
Google\textquotesingle s APIs, not a passing attestation: microG obtains
at most the software basic-integrity verdict, while the hardware-backed
device and strong verdicts require a certified device with a locked
bootloader and stay out of reach \citep{eos-microg}, so the banking- and
payment-application refusals of Section 8 persist on these systems -\/-
CalyxOS\textquotesingle s own guidance states that compatibility depends
on the specific app and that apps requiring full Play Services
integration, such as Google Pay, will still not work
\citep{calyxos-banking}. That spoofing path is itself narrow and
shrinking: as Section 8 describes, only the basic and legacy
pre-Android-13 verdicts fall to software spoofing -\/- Universal
SafetyNet Fix states that "Android versions up to 13 are supported" and
has been unmaintained since 2023 \citep{safetynet-fix} -\/- while the
hardware-backed verdict on devices launched with Android 13 or later
yields only to a leaked attestation keybox, and those are not durable,
publicly shared keys being banned by Google as they surface
\citep{playintegrityfork}, \citep{tricky-store}; it is the arms race
Section 14 describes, in which the spoofing side has been progressively
constrained \citep{nguyen-vu-android-rooting}.

That bind raises a further question -\/- whether any owner-controlled
privilege could escape it rather than merely confirm it, a verified
owner-inspection mode that preserves both verified boot and attestation
-\/- which Section 13 takes up, systematizing the design space and the
gap none of its approaches closes.

Aftermarket operating systems are thus consistent with the
paper\textquotesingle s analysis: owner-inspectability is attainable,
but only by leaving the as-sold, vendor-keyed configuration behind,
after which the attestation ecosystem imposes the functional penalty
described in Section 8. For the overwhelming majority of devices, which
run their as-sold vendor software, the inspection limits catalogued here
are unchanged.

One limit even this configuration cannot reach is that the owner still
cannot independently image or re-verify every component the device runs.
The cellular baseband is the clearest case: it executes its own
vendor-controlled firmware on a separate processor that the
application-processor operating system -\/- stock or aftermarket -\/-
does not inspect, and remotely exploitable memory corruption in a
baseband stack is a demonstrated class of attack
\citep{weinmann-baseband}. On several Samsung Galaxy models the
proprietary radio-interface-layer software on the application processor
-\/- the vendor library that mediates the modem -\/- was moreover found
to implement a remote-file-service protocol through which the modem can
read, write, and delete files on that storage; its finders characterized
this as a backdoor, though whether it was a deliberate backdoor rather
than an incidental feature of the modem\textquotesingle s file service
is not established \citep{replicant-samsung-backdoor}. More recent
analysis of Samsung devices continues to uncover undisclosed baseband
command interfaces and vendor-confirmed zero-day flaws, including denial
of cellular service and arbitrary access to the application
processor\textquotesingle s file system from the baseband
\citep{li-basemirror}. The Samsung finding was itself obtained by
reverse-engineering that proprietary vendor library, a reminder that
realizing inspectability in practice can draw on the
modify-and-reverse-engineer freedoms held conceptually distinct from it
in Section 3, and can carry legal exposure the owner or researcher must
navigate: the finders rested the lawfulness of their decompilation on
the interoperability exception in Article 6 of the EU Computer Programs
Directive \citep{replicant-samsung-backdoor-wiki},
\citep{eu-computer-programs-directive}.

A discrete secure element, where present, likewise runs vendor firmware
on its own chip. The trusted execution environment should not be
conflated with these: on a platform such as this one it shares the
application processor and is signed and checked by the bootloader, so it
falls under verified boot rather than being an uninspected separate
processor; its limitation for the owner is that its firmware, like the
baseband\textquotesingle s, is not something the owner can independently
image and audit, not that it runs unverified. The broader pattern is
that persistent implants have migrated toward low-level firmware beneath
the operating system, installable by ordinary malware once the system is
compromised \citep{zaddach-hdd-backdoor}; the strongest
owner-inspectable configuration still cannot independently audit that
firmware layer, echoing in a new place the principle that a system
cannot be trusted to report truthfully on its own state
\citep{thompson-trusting-trust}.

A third objection holds that verified boot and hardware key attestation
already furnish the integrity assurance that owner-inspection is meant
to provide, so the owner has no need to verify the device independently.
Verified boot is itself a strong defense but not an absolute one, and
the owner cannot independently confirm that it holds. Its guarantees
rest on assumptions -\/- that the signing keys stay secret, and that the
bootloader and trusted execution environment enforcing the chain are
themselves correct -\/- and these have been broken. Platform signing
keys have leaked and been used to sign malware in the wild
\citep{apvi-platform-certs}, and bootloader and
trusted-execution-environment compromises that would subvert the chain,
including extraction of the hardware key material that underpins
attestation, have been demonstrated in research and are reachable by
malware that already controls the operating system, though none has been
observed as an in-the-wild boot implant on a locked device
\citep{bootstomp}, \citep{shakevsky-trustzone}. Such malware cannot
simply overwrite the bootloader -\/- the immutable hardware root of
trust above it verifies the bootloader\textquotesingle s signature and
refuses a tampered image -\/- but it can exploit a flaw in the signed
bootloader or trusted code to execute within the chain, which is how a
software-only compromise reaches beneath the operating system.

The hardware root of trust that anchors this check is itself a physical
assumption rather than an absolute: the boot ROM is typically immutable
mask ROM, so where it is a flaw in it is permanent and unpatchable, and
glitching it can bypass signature verification altogether
\citep{ncc-mediatek-bootrom} -\/- a class not specific to one vendor
\citep{checkm8-vu}; some systems-on-chip instead expose a ROM-patch
path: a small vendor-signed patch, held in one-time-programmable fuses,
that redirects the boot ROM\textquotesingle s early execution and is
itself locked by a fuse \citep{opentitan-rom-patch}. Immutability is
therefore a vendor design choice rather than a constant, and either way
the mechanism is the vendor\textquotesingle s and not the
owner\textquotesingle s; and being signed code gated by fuses, that
patch path is itself subject to the same key-compromise and glitching
failure modes, and adds early-boot attack surface of its own. An
invasive silicon attack can further read out the fused device keys
beneath the chain \citep{ioactive-antifuse}; but every such defeat
documented to date requires physical possession of the device and
belongs to the forensic and unlock class, not to any remote, in-the-wild
implant on a locked consumer phone. A remote-seeded, locally-triggered
variant -\/- malware or a malicious update planting malformed data that
an early boot stage parses on the next boot -\/- is conceivable in
principle, but none has been documented, and on a locked device the code
the boot ROM loads is signature-verified against a fuse-anchored key
\citep{aosp-verified-boot}, so a purely over-the-air boot-ROM compromise
is not a realistic or observed attack surface.

The deeper point for the owner is epistemic: on the stock, vendor-keyed
configuration the green boot state the device shows and the attestation
certificate it presents are the platform\textquotesingle s own
self-report, which a sufficiently deep breach can forge and which the
owner has no independent means to check. This limit is the hardware
beneath the chain, not the verified-boot signing key above it: on
devices that support it, Pixel among them, the owner can enroll a custom
signing key and re-lock to it, reaching a yellow boot state whose
on-screen key fingerprint the owner can themselves verify
\citep{aosp-boot-flow}, \citep{aosp-device-state} -\/- what stays beyond
any owner is the boot ROM and the fused keys under that key, which none
can replace or read out.

A fourth objection is that consumer antivirus and mobile-security apps
-\/- products such as Kaspersky, Avast, and Norton, alongside
Google\textquotesingle s built-in Play Protect -\/- already let an owner
detect compromise. These are real and widely used: Play Protect is, in
Google\textquotesingle s description, "the most widely deployed mobile
threat protection service in the world" and "scans 200 billion Android
apps daily" \citep{google-play-protect}. But all of them, however
privileged, run inside the very system they are meant to vet. The
third-party products operate as ordinary applications, while Play
Protect is more privileged still -\/- a component of Google Mobile
Services that reaches across the device to disable or remove other apps
on the owner\textquotesingle s behalf
\citep{google-play-protect-remove}, \citep{google-gms} -\/- yet it too
executes within the running operating system rather than from any
external, trusted vantage. They match applications and files against
known-malware signatures and heuristics within the storage the operating
system permits them to see -\/- Kaspersky\textquotesingle s full scan,
for instance, covers "the device\textquotesingle s entire file system"
\citep{kaspersky-android} -\/- which catches known, application-level
threats but presupposes an operating system that is reporting truthfully
about itself.

The vendors say as much. Avast states the limit plainly: "Like all
applications on the Android market, it is subject to limitations imposed
by the operating system version running on the device. It cannot protect
against exploits targeting specific vulnerabilities in the OS kernel,
the network stack and other fundamental parts of the system."
\citep{avast-mobile-faq} That is exactly the class of compromise
owner-inspectability concerns: a kernel-level or sub-operating-system
implant sits beneath the vantage point these tools occupy and can
neutralize or evade an in-system scanner whatever its privileges, as
Sections 9 and 14 develop. Consumer antivirus therefore catches the
known and the application-level but does not supply the external,
trusted vantage point verification requires -\/- narrowing the gap for
commodity malware without closing it for the targeted, low-level
compromise, and so consistent with the analysis above.

\subsection{13. Owner-Controlled Inspection and the Verified-Inspection
Gap}\label{owner-controlled-inspection-and-the-verified-inspection-gap}

The objections above return repeatedly to a single missing capability.
Owner-inspectability is foreclosed because the privileged access an
owner needs in order to inspect the system -\/- to examine it from a
clean vantage, to detect a kernel-level rootkit or vendor-supplied code
the owner cannot audit, and to let independent researchers work on
hardware they own -\/- is exactly the access that forfeits the
device\textquotesingle s attested standing. The capability this bind
calls for is what this paper calls a verified owner-inspection mode:
owner-controlled privileged inspection that surrenders neither verified
boot nor the attestation a relying party accepts. This section asks
whether such a capability exists. It proposes and evaluates no system;
it systematizes the approaches that have been shipped or proposed,
locates each against the verified-boot security model, and identifies
the gap none of them closes -\/- an open problem whose resolution would
widen owner-inspectability well beyond the few who can tolerate the
attestation penalty.

Several security-focused aftermarket systems decline to support root for
this reason. CalyxOS "does not support running in rooted mode, nor do we
recommend this," on the ground that "Running any Android device with
root permissions severely undermines the security of the device"
\citep{calyxos-root-faq}; the discontinued DivestOS likewise "does not
support or encourage the use of root or runtime modification
frameworks," which "break the trust model that AOSP employs and reduce
the integrity of the system" \citep{divestos-root-faq}.

That trust model is itself platform-centric: the Android platform
security model requires agreement among user, platform, and developer,
and it classes rooting, "both intentional and malicious," as "a
non-compliant change which violates CDD" -\/- the Compatibility
Definition Document a device must satisfy to ship as Android
\citep{aosp-cdd} -\/- holding that ownership is "a legal concept rather
than a technical one and therefore outside the scope of an OS security
model" \citep{mayrhofer-android-security-model}.

GrapheneOS has weighed the question directly, and its lead developer has
addressed it in more technical detail: a carefully constrained root -\/-
"an on-device root terminal which requires a special boot mode and
isn\textquotesingle t persistent," limited to the user and barred from
hijacking by apps -\/- is "theoretically possible to do it in a way that
doesn\textquotesingle t throw away a lot of security," but
"isn\textquotesingle t being done by anyone right now"; the phrase
"totally ruining verified boot by trusting persistent state" is directed
instead at the persistent, fully-trusted way the root-enabled forks
integrate it \citep{grapheneos-root-statement}. The
project\textquotesingle s own account put that cost plainly: once
persistent state and the application layer are fully trusted, "the
security model for verified boot and attestation are not intact to
detect" a resulting compromise \citep{grapheneos-root-forum}.

The boundedness of root established in Section 3 is thus not merely a
restriction the stock vendor imposes and an aftermarket system could
lift for free: a persistent root exacts a cost, and that cost is borne
by the owner who enables it. An owner can relock under a self-signed key
and keep cryptographic verified boot even with root present
\citep{resign-android-image}, and the vendor-anchored attestation
verdict that stock relying parties demand is in any case already
forfeited by leaving the vendor\textquotesingle s key behind, rooted or
not; but a persistent, fully-trusted superuser defeats the verified-boot
security model of not trusting persistent state, so the self-signed
state the owner could otherwise check no longer tells them the device is
untampered. That owner-side verifiability is what a persistent root
costs, and a root that keeps it is, by these projects\textquotesingle{}
own accounts, unbuilt where it is thought possible at all and rejected
outright elsewhere. Hardware attestation is designed around exactly
this: GrapheneOS\textquotesingle s own documentation notes that the
attestation API reports the device\textquotesingle s state "in a way
that even an attacker exploiting the OS after boot to gain root cannot
trivially bypass" \citep{grapheneos-attestation-compat}.

Restricting attention to what an owner can do while keeping the
bootloader locked and verified boot engaged, the approaches differ on
what actually varies among them (Table 4): whether they grant a full
superuser, and whether they preserve the verified-boot security model
under which the owner can still detect tampering. The pattern is that a
full root forfeits that security model even on a locked, owner-signed
device -\/- a persistent, fully-trusted superuser is the compromise the
model exists to catch -\/- while the configurations that keep it are
either non-persistent or not a full superuser at all.

Table 4. Owner root and privilege options that keep the bootloader
locked and verified boot engaged.

\begin{longtable}[]{@{}
  >{\raggedright\arraybackslash}p{(\columnwidth - 4\tabcolsep) * \real{0.3333}}
  >{\raggedright\arraybackslash}p{(\columnwidth - 4\tabcolsep) * \real{0.3333}}
  >{\raggedright\arraybackslash}p{(\columnwidth - 4\tabcolsep) * \real{0.3333}}@{}}
\toprule\noalign{}
\begin{minipage}[b]{\linewidth}\raggedright
Approach
\end{minipage} & \begin{minipage}[b]{\linewidth}\raggedright
Full root?
\end{minipage} & \begin{minipage}[b]{\linewidth}\raggedright
Verified-boot security model preserved (owner can detect tampering)?
\end{minipage} \\
\midrule\noalign{}
\endhead
\bottomrule\noalign{}
\endlastfoot
Root baked into a signed image and relocked under an owner key: a
Magisk-style systemless patch (avbroot \citep{avbroot}) & Yes & No -\/-
verified boot no longer certifies the running system: the persistent,
fully-trusted superuser and its module store (in /data/adb/modules
\citep{magisk-guides}) live above the verified layer, in unverified
/data \citep{aosp-dm-verity}, \citep{aosp-file-based-encryption} \\
Scoped kernel root baked in the same way: KernelSU, with per-app
profiles \citep{kernelsu} & Yes, scoped per app & No -\/- a persistent
kernel-level superuser; app-profile scoping narrows the surface but is
not a security boundary \\
Non-persistent root: a special-boot-mode terminal not retained across
reboot \citep{grapheneos-root-statement} & Only transiently & Yes -\/-
the root does not persist across reboot, which restores the verified
boot chain and system image (writable user data, outside verified boot,
is a separate matter) \\
Signed, verified extension shipped inside the image: Android APEX
\citep{aosp-apex} (the Linux analogue, under a verity-and-signature
image policy, is systemd-sysext \citep{systemd-sysext}) & No -\/-
bounded, policy-confined verified capability & Yes -\/- a signed,
verity-checked image, so its contents are verified rather than arbitrary
mutable state \\
ADB-root baked into an owner-signed, relocked image: adbd patched to run
as root, with no on-device su (resign-android-image
\citep{resign-android-image}) & No -\/- root is reachable only over the
ADB interface; no app-facing superuser & Yes -\/- the root path ships
inside the owner-signed, verity-covered image, so tampering remains
detectable and a reboot restores the verified state; it trusts no
unverified persistent state, unlike a /data module store \\
Privileged but not root: ADB-shell brokering (Shizuku) \citep{shizuku} &
No -\/- capped at the shell permission set & Yes -\/- stock image
unmodified \\
\end{longtable}

Three qualifications attach to Table 4. The non-persistent row can be
realized by a RAM-backed overlay, but an overlay\textquotesingle s upper
layer lies outside the verified-boot hashtree \citep{aosp-dm-verity}, so
the content it serves is unverified whatever backs it; ephemerality
bounds the window of unverified operation rather than keeping the
running system verified. And the design space is not closed in
principle: image-based operating systems have prompted proposals for an
owner-facing "developer mode" that would reconcile locally built or
signed code with an immutable, cryptographically validated system, but
such a mode remains a design proposal rather than a shipping feature
\citep{poettering-developer-mode}. Third, the first row assumes the
verity-preserving case: avbroot and resign-android-image keep dm-verity
enabled under the owner\textquotesingle s key, so cryptographic verified
boot stays intact and only its security model is lost; a Magisk-only
patch that disables dm-verity instead abandons even cryptographic
verified boot.

A concrete form illustrates the persistent-root row. An owner could bake
a root-accessible SSH daemon into the owner-signed image and authorize
it with the very key that signs that image -\/- the
owner\textquotesingle s verified-boot public key placed in the
daemon\textquotesingle s list of authorized keys, so that a single owner
key governs both boot verification and administrative login while its
private half stays off the device. Cryptographic verified boot is
preserved and the key hierarchy is consistent, yet a persistent root
service is exactly the fully-trusted persistent state the model is built
to distrust -\/- a verdict that holds whether it is reached locally or
over the network, since what the model declines to trust is the
persistent, fully-privileged state itself, not the transport that
reaches it. A key-only daemon is in itself a mature and widely accepted
surface; what the network transport adds is narrower -\/- compromise no
longer requires physical possession of the device, which leaves behind
the on-device, physically-present case on which the safe reading
depends. Gated instead to a non-persistent maintenance mode, it
collapses into the non-persistent row.

Role separation sharpens the same conclusion. Confining administrative
capability to a separate account or a deliberately entered maintenance
session, while the owner\textquotesingle s daily session runs
unprivileged with no reachable escalation path, applies the principle of
least privilege \citep{saltzer-schroeder} and separation of duty
\citep{clark-wilson} and keeps the everyday attack surface off the
privileged path. But the account boundary alone does not preserve the
verified-boot model: a persistent, writable administrative role is
offline persistent root, so the model survives only if that maintenance
context is itself ephemeral or re-sealed -\/- its integrity re-measured
against the owner\textquotesingle s own key -\/- on exit. In the
owner-signed configuration the owner does hold that key and can re-seal
\citep{avbroot}; the constraint is operational: secure practice keeps
the private half off the device (only the public key is enrolled), so
re-sealing is an off-device re-sign and reflash from a trusted
workstation rather than an in-situ maintenance step, and keeping the key
on the device to make it in-situ would reintroduce the persistent
trusted secret the locked chain exists to avoid. Such a re-seal restores
the owner\textquotesingle s own verified state, not the vendor-anchored
attestation verdict, and no shipping system integrates it.

No approach in this design space, then, delivers all three of persistent
owner administrative control, the verified-boot security model, and a
relying-party-accepted attestation at once. The configurations that keep
the security model do so by giving up persistence -\/- the
non-persistent root -\/- or full privilege -\/- a verified extension, or
shell-tier brokering -\/- while the one that keeps full persistent
privilege under the owner\textquotesingle s own key gives up the model
and, with its non-vendor key, relying-party attestation. A capability
that closed the gap would have to assemble three things the shipped
approaches reach only in pairs: privileged inspection the owner can
invoke; an integrity state the owner, and ideally an independent
verifier, can still check afterward, which points to non-persistent or
verified-extension privilege with an owner-held re-seal, subject to the
off-device-key constraint above; and, for the relying-party half, an
independent trust anchor for owner-enrolled keys that the attestation
ecosystem does not presently provide. These three requirements are not
alike in who can satisfy them. The first two are device-level and within
the owner\textquotesingle s reach: non-persistent or verified-extension
privilege under an owner-held key already assembles privileged
inspection with an integrity state the owner, and an independent
verifier, can re-check, so the owner-side half is largely a matter of
combining shipped mechanisms rather than inventing new ones. The third
is not the owner\textquotesingle s to satisfy. A relying-party-accepted
verdict is anchored, by construction, to the vendor\textquotesingle s
root of trust, and the attestation ecosystem does not accept
owner-enrolled keys (Appendix A); nor, on the available evidence, does
it have reason to, since the same verified, vendor-keyed state is a
precondition for hardware-backed content protection
\citep{patat-widevine}, \citep{rafi-drm-mobile}, \citep{aosp-drm} and
for the platform-policy objectives the model is built to enforce
\citep{mayrhofer-android-security-model},
\citep{doctorow-general-computation}. The open problem is therefore
asymmetric. Owner-side verifiability is attainable today for owners who
accept the loss of vendor-gated attestation; re-admitting an
owner-controlled device to the relying-party ecosystem is not a missing
device mechanism but a policy choice that rests with the platform, and
the content-protection and platform-control interests give it a standing
reason to withhold. Treating the owner-side half as an open engineering
problem, rather than a settled impossibility, is what would let
owner-inspectability -\/- and with it independent malware and backdoor
investigation -\/- reach beyond the owners who can accept the
attestation penalty.

\subsection{14. Discussion}\label{discussion}

The limitations examined here are not independent defects but facets of
a single design choice: administrative capability is withheld from the
device owner and conferred only on the vendor and manufacturer. Verified
boot, sandboxing, mandatory access control, attestation, and
hardware-bound key storage each advance the platform\textquotesingle s
assurance that only approved, unmodified software runs, and each
simultaneously removes a means by which an owner might inspect the
device from a trusted vantage point. The circular trust dependency of
Section 8 is the sharpest expression of the tension: the configurations
an owner would need in order to verify the device are exactly those the
platform refuses to attest as trusted.

The owner is consequently placed in a position analogous to that of an
attacker with respect to introspection, and in one respect a more
constrained one: these controls foreclose the owner\textquotesingle s
sole avenue -\/- modifying the device to inspect it -\/- while leaving
open the runtime and sub-operating-system routes the capable adversary
of Sections 8 and 12 uses, and the owner brings none of an
attacker\textquotesingle s willingness to damage the device. The
practical recourses -\/- JTAG-based memory acquisition, chip-level
storage extraction, and exploit-driven temporary root -\/- are,
respectively, rarely available, destructive or specialized, and
dependent on the existence of exploitable vulnerabilities that the
platform is otherwise trying to eliminate. The physical-acquisition path
is further foreclosed where unlocking the bootloader triggers a
mandatory wipe of user data, so that gaining access to the storage also
erases its contents \citep{aosp-bootloader-locking}. None of these
recourses offers the routine, non-destructive, owner-driven verification
that the threat model of Section 3 calls for, and, as Sections 3 and 8
established, even a successful temporary root yields only a bounded
in-system privilege on a device that has already forfeited its attested
standing -\/- not the external, trusted vantage point that verification
requires. Nor would more privilege close the gap: a kernel-space rootkit
manipulates the very kernel structures that any in-system observer, root
included, must query, so it can neutralize a privileged security tool
and falsify its results regardless of the privilege level at which that
tool runs \citep{guri-joker-jtag},
\citep{bickford-rootkits-smartphones}. Subversion-resistant detection
must therefore observe the system from outside the running operating
system \citep{guri-joker-jtag}, \citep{garfinkel-vmi}, echoing the
long-established principle that a system cannot be trusted to report
truthfully on its own state \citep{thompson-trusting-trust}.

The contrast with a fully controlled host is one of category, not merely
degree. On a desktop or server the baseline acquisition workflow -\/-
detach the storage and image it through a write blocker, a device that
passes read commands while blocking any that would alter the drive, so
the acquisition does not modify the evidence \citep{nist-sp800-101r1},
\citep{nist-hwb-spec}, then decrypt the image with the
owner\textquotesingle s passphrase-derived key -\/- is available to an
owner who can remove the storage and holds the decryption key, and
remains so even though the storage device itself runs opaque controller
firmware that a determined attacker can subvert: compromising the
firmware of a commercial off-the-shelf drive to backdoor data as it is
written has been demonstrated using only public information and reverse
engineering, installable by ordinary malware once the operating system
is compromised \citep{zaddach-hdd-backdoor}, and such drive-firmware
implants have been found in the wild \citep{kaspersky-equation}, with
persistence that survives disk formatting and a full operating-system
reinstall \citep{kaspersky-equation-qa}. Implants one layer lower still,
in the motherboard\textquotesingle s own UEFI flash, persist even across
a storage replacement, because they do not reside on the storage at all
\citep{eset-lojax}. That residual is real, but it is a different thing
from the phone\textquotesingle s barrier: an advanced and partially
mitigable attack vector layered beneath an acquisition path that
otherwise works, not a categorical barrier for all owners.

On a stock phone the same baseline workflow is foreclosed at its first
step, because the storage is soldered, and at its last, because even a
successful chip-level read returns the user\textquotesingle s data as
ciphertext under a hardware-bound key. The firmware-implant concern,
common to both platforms and marking the floor of what either owner can
verify, thus compounds on the phone a barrier that is already
categorical rather than standing, as on the host, as the one residual
beneath an otherwise working path.

Hardware memory acquisition sharpens the same contrast, but only partly,
because it is closing on both sides. On a conventional computer an
examiner can in principle capture RAM from outside the running operating
system -\/- by direct memory access over an expansion bus
\citep{markettos-thunderclap}, \citep{frisk-pcileech} or by exploiting
the residual charge in DRAM after a cold reboot
\citep{halderman-coldboot} -\/- yet modern platforms increasingly block
such access with an input-output memory management unit
\citep{markettos-thunderclap}, and cold-boot capture needs physical
access to the memory. A stock phone forecloses even this residual
avenue: its memory is soldered in package-on-package stacks that cannot
be transplanted, it exposes no external bus capable of such access, and
the one demonstrated on-device cold-boot variant first required
unlocking the bootloader, which wipes the user data it sought to recover
\citep{muller-frost-coldboot}.

Table 5 sets the owner\textquotesingle s control over each layer of the
stack side by side for the two platforms, read downward from the
bootloader to the firmware beneath the operating system. The pattern is
a categorical difference across the upper, operating-system-and-boot
layers, where the personal-computer owner holds administrative control
and the stock-phone owner holds none, and a shared opaque floor at the
firmware and management-coprocessor layers below them, where neither
owner can inspect what runs.

Table 5. Owner control over the stack: a traditional modular personal
computer the owner administers versus stock Android (as sold).

\begin{longtable}[]{@{}
  >{\raggedright\arraybackslash}p{(\columnwidth - 4\tabcolsep) * \real{0.3333}}
  >{\raggedright\arraybackslash}p{(\columnwidth - 4\tabcolsep) * \real{0.3333}}
  >{\raggedright\arraybackslash}p{(\columnwidth - 4\tabcolsep) * \real{0.3333}}@{}}
\toprule\noalign{}
\begin{minipage}[b]{\linewidth}\raggedright
Stack layer
\end{minipage} & \begin{minipage}[b]{\linewidth}\raggedright
Traditional personal computer (owner)
\end{minipage} & \begin{minipage}[b]{\linewidth}\raggedright
Stock Android, as sold (owner)
\end{minipage} \\
\midrule\noalign{}
\endhead
\bottomrule\noalign{}
\endlastfoot
Bootloader & Replaceable; Secure Boot is optional, and the owner can
enroll custom keys or disable it & Locked and vendor-signed; unlocking
wipes user data and drops attestation, and some devices are unlockable
only by exploit or not at all \\
Kernel & Owner builds and boots an arbitrary kernel & Vendor-signed boot
image; a custom kernel requires unlocking \\
Kernel modules & Owner loads arbitrary modules & Lockdown,
module-signature, and dm-verity enforcement restrict loading \\
initramfs & Owner controls & Part of the signed boot image; not
owner-modifiable on stock \\
init (first userspace, PID 1) & Owner chooses or replaces the init
system, and can boot to an alternate init or single-user mode & Android
init ships in the signed ramdisk and system partitions under verified
boot and dm-verity, runs under SELinux, and is not owner-modifiable on
stock \\
Operating system & Owner installs any operating system & A custom
operating system requires unlocking and incurs an attestation penalty \\
In-system superuser (root) & Available to the owner by default &
Withheld as sold; reachable only by exploit or unlocking -\/-
exploit-gained root is SELinux-confined, while systemless root sidesteps
SELinux but remains below the kernel\textquotesingle s lockdown and
module-signature enforcement \\
Boot an external trusted system & Yes, from external media & Not
supported while locked: the bootloader boots only signed images and an
external system requires first unlocking \citep{aosp-fastboot-protocol},
\citep{aosp-bootloader-locking}; the path is then device-specific \\
Remove and image storage & Yes: detach the storage and image it on a
trusted controller & Storage is soldered; chip-level reading is
destructive, and the data returns as ciphertext under a hardware-bound
key \\
Platform firmware & BIOS/UEFI: reflashable but vendor-signed and opaque
in practice, and in-the-wild implants have been demonstrated
\citep{eset-lojax} & Bootloader and vendor firmware: vendor-controlled
and owner-opaque \\
Management co-processor & Intel ME or AMD PSP: a vendor-firmware
co-processor the owner cannot inspect -\/- the ME embedded in the
chipset with privileged access to platform data
\citep{ermolov-intel-me}, the AMD-SP a boot-time root of trust that runs
before the main cores \citep{jacob-faultpm} & Cellular baseband, and any
discrete secure element: vendor firmware on separate silicon the owner
cannot inspect \citep{weinmann-baseband} \\
Storage-controller firmware & Opaque; controller-firmware implants
demonstrated in the lab \citep{zaddach-hdd-backdoor} and found in the
wild \citep{kaspersky-equation} & Opaque \\
\end{longtable}

The division is systematic. On the personal computer the owner commands
the entire operating-system-and-boot stack and is blocked only at the
firmware and management-coprocessor floor; on the stock phone the owner
commands essentially none of the stack, so the same firmware floor sits
beneath a barrier that is already categorical at every layer above it.
Administrative control, in the sense defined in Section 3, is exactly
what the left column has and the right column lacks; root, where a phone
owner reaches it, adds only the single "in-system superuser" row and
leaves every other row unchanged.

The ease of this recourse may be overstated. Rooting presumes an act
that retail devices are configured to resist: they ship with the
bootloader locked, and unlocking it -\/- where the vendor permits it at
all -\/- is gated behind an explicit warning and a mandatory wipe of
user data \citep{aosp-bootloader-locking}, \citep{aosp-boot-flow}. Many
carrier- and vendor-locked devices cannot be unlocked by their owners at
all, or only by means of an exploit \citep{bootloader-wall-of-shame}.
Unlocking is widely said to void the device warranty, although
consumer-protection law in the United States and the European Union
limits blanket voiding and requires a defect to be causally linked to
the modification before coverage may be refused
\citep{magnuson-moss-act}, \citep{eu-directive-2019-771}. Several
mainstream vendors attach irreversible hardware-level consequences to
unlocking, by distinct mechanisms rather than a single universal one:
one permanently trips a one-time-programmable e-fuse that disables
hardware-backed attestation and the protected features that depend on it
and cannot be reset \citep{samsung-knox-hw-security}; another
irrevocably erases device-unique keys that relocking does not restore
\citep{sony-unlock-bootloader}; and one-time-programmable (eFuse)
hardware is, more generally, the standard system-on-chip primitive for
secure-boot enablement and anti-rollback enforcement, whose blowing is a
one-way operation \citep{qualcomm-secure-boot}.

A countervailing point is that not every platform is irreversible: on
some devices unlocking can be reversed and the verified-boot and
attestation state is restored on relocking -\/- to the full state under
the vendor\textquotesingle s stock key, or to a reduced state under an
owner-enrolled key \citep{aosp-verified-boot},
\citep{aosp-key-id-attestation}. Even there the bind holds, because the
reversible path returns the owner to the vendor-keyed configuration
without ever granting the top attestation standing while the device
remains open to inspection.

Taken together, these frictions make unlocking a high-barrier step
rather than a routine one. The privileged state an owner does reach is,
as Section 8 established, counterproductive for verification; and the
integrity-masking tools owners deploy to pass integrity checks work by
presenting the profile of an unmodified stock device, leaving them in an
arms race against detection in which the spoofing side has been
progressively constrained \citep{nguyen-vu-android-rooting}.

The unlock procedure is, moreover, not silent to software already
resident on the device. Preparing a device for inspection changes
observable system state, and unlocking the bootloader triggers a
mandatory wipe that destroys the user-data partition before any image of
it can be taken -\/- so the inspection a high-threat owner attempts can
itself alert, and then erase, a resident implant.

That such environmental signals drive evasion is established rather than
speculative: measurement studies show modern malware checks for
virtualized or instrumented environments and withholds its behavior when
it detects them \citep{chen-anti-virtualization},
\citep{ruggia-evasive-malware}, and the anti-forensics literature names
the matching capabilities of artifact wiping, trail obfuscation, and
self-removal \citep{garfinkel-anti-forensics},
\citep{mitre-attack-evasion}. Mobile spyware has combined the two in the
wild: NSO\textquotesingle s Pegasus for Android shipped "suicide
functionality to remove itself" \citep{lookout-chrysaor} (Suicide
Functionality section), triggered among other cases when it cannot find
a valid cellular subscriber identity -\/- behavior attributed to
avoiding analysis and emulator environments -\/- while the iOS variant
self-destructs when "a great probability of exposing the agent exists"
\citep{marczak-million-dollar-dissident} (Section 3.5), and later
versions were observed deleting their own process records to hide their
traces \citep{amnesty-pegasus-methodology}.

The specific loop -\/- an implant watching for unlock preparation and
erasing itself before acquisition -\/- remains a reasoned extrapolation
rather than a documented exploit; its constituent capabilities are each
demonstrated, and the principle that a compromised system cannot be
trusted to report truthfully on itself is long established
\citep{thompson-trusting-trust}, \citep{garfinkel-vmi}. The nearest
realized case, the Rombertik malware, destroys its host on detecting
that its own code has been analyzed \citep{talos-rombertik}; this is a
genuine detect-then-destroy reaction, but it is keyed to tampering with
the malware\textquotesingle s own image rather than to a forensic
acquisition, and it wrecks the host rather than quietly erasing the
implant so the device passes inspection clean. No reported case closes
the precise loop posited here.

\subsubsection{Limitations of this
analysis}\label{limitations-of-this-analysis}

Several limitations bound the claims made here. The paper is a
systematization of existing evidence, not a new empirical study: it
reports no measurements of its own, and where the literature is thin
-\/- notably on the precise prevalence of attestation-gated applications
and on the fraction of devices for which hardware acquisition is
feasible -\/- it says so rather than substituting an estimate. Its
central construct, owner-inspectability, is a qualitative criterion
rather than a metric; this paper defines and applies it but does not
quantify it, and an operational scale by which devices and
configurations could be scored is left to future work. The analysis is
deliberately scoped to the owner\textquotesingle s standpoint and to the
Android platform as shipped on mainstream devices; it does not weigh the
owner-side costs against the genuine ecosystem-wide benefits that the
same mechanisms deliver against running adversaries, which the platform
literature documents \citep{mayrhofer-android-security-model}, and it
does not extend its specific claims to other platforms except where a
cited source does so. Finally, the platform is a moving target: verdict
tiers, acquisition techniques, and project capabilities change between
releases, and the specific mechanisms cited here are current as of
writing rather than permanent.

Inspectability and modifiability are distinct as ends -\/- the criterion
asks that the owner establish and demonstrate what the device runs, not
that the device be freely modifiable -\/- yet on locked stock devices
they remain entangled in practice, in three ways. Mechanically, gaining
a first inspection foothold on such a device generally requires the same
unlock that also admits an arbitrary replacement system, so the
capability the owner needs tends to arrive bundled with one the platform
withholds. Through the software-freedom ecosystem, the freedoms to
study, rebuild, and redistribute sustain inspection from upstream:
auditable source and reproducible builds give the owner a known-good
reference to check an acquired image against -\/- they fix what a
trusted build should contain, though confirming the device actually runs
it still takes an acquired image or a trusted measurement -\/- and
shared tools and findings are what make inspection practical at all. And
since inspectability is only a proxy for the owner\textquotesingle s
actual end -\/- security for the owner, not the platform -\/-
modification can reach that end by another route: replacing an opaque
firmware component with a more trustworthy build can improve security
where inspecting it is impossible, though it substitutes a better-placed
trust rather than delivering the verification the criterion measures,
the same move as re-locking to a different vendor\textquotesingle s key
(Section 12). That reaching inspection on a locked stock device usually
routes through the broader modification the platform withholds is itself
the case for decoupling the two by design.

A measured reading is that the platform optimizes for resistance to the
ordinary running adversary at the cost of owner-side auditability, and
that the two objectives, while not strictly incompatible, are in current
designs rarely satisfied together. The forensic literature surveyed here
does not claim that compromise is undetectable on stock devices; it
claims that detection is constrained by observability, and that a
negative result is accordingly weak evidence of integrity. Building on
that, and on the user-controlled root of trust that AOSP allows in
principle (Section 3) and the owner-signed GrapheneOS configuration of
Section 12, this paper\textquotesingle s own position is that
auditability could be improved without abandoning existing protections.
That position suggests a design space -\/- owner-attestable inspection
modes, exportable integrity measurements, or supported trusted-boot
verification paths -\/- in which verifiability and platform integrity
might be advanced together rather than traded off.

\subsection{15. Conclusion}\label{conclusion}

Vendor-locked devices systematically limit the owner\textquotesingle s
ability to inspect, control, and independently verify the software they
run. Stock Android, among the most openly specified such platforms and
this paper\textquotesingle s worked example, is the case developed here:
the owner cannot practically examine the whole system from a clean,
trusted environment; application and update delivery is mediated by
centralized providers whose abuse could enable targeted malicious
installation; internal storage and live memory are inaccessible without
withheld privileges; attestation couples a device\textquotesingle s
attested standing to the device remaining unmodified; low-level rootkit
detection generally demands a trusted memory snapshot taken from outside
the running system; and the established techniques of host-based malware
analysis are, as a class, unavailable without first breaking the
verified-boot and attestation guarantees the platform enforces.

Obtaining root narrows these limits without lifting them. It lets an
owner inspect at-rest storage, installed packages, and process memory
-\/- enough to find commodity, application-level, and at-rest malware
-\/- but not the targeted low-level class: root remains an in-system
vantage point that a kernel-space rootkit can deceive, it does not reach
the kernel, boot chain, or firmware beneath it, and -\/- whether gained
by exploit or by unlocking -\/- it forfeits the attested status the
owner is trying to preserve. Backing up the hardware-bound encryption
key is likewise impractical for most owners.

The existence of consensual forensic tooling, and of aftermarket
operating systems such as GrapheneOS, does not overturn these
conclusions: such tooling is bounded by the limited diagnostic data the
platform exposes, owner-inspectable configurations are reachable only by
leaving the as-sold device behind and paying an attestation penalty -\/-
the loss of device-integrity standing that leads an apparently growing
set of applications to refuse to run -\/- to do so, and an absence of
findings is accordingly weak evidence of integrity. These limits are not
accidents of one vendor\textquotesingle s engineering but consequences
of vendor-lock itself, and the owner\textquotesingle s structural
position recurs, to varying degrees, on other vendor-locked devices.
Treating security as owner-side verifiability, rather than as
platform-side policy enforcement, makes these limitations visible as a
coherent whole and identifies owner-inspectability as a property in need
of improvement.

\subsection{Appendix A. GrapheneOS Auditor verified-boot key
allowlist}\label{appendix-a.-grapheneos-auditor-verified-boot-key-allowlist}

Section 12\textquotesingle s claim that the released GrapheneOS Auditor
recognizes only keys compiled into it rests on the app\textquotesingle s
own source. Accepted verified-boot keys are fixed
\texttt{ImmutableMap\textless{}String,\ DeviceInfo\textgreater{}} tables
keyed by the hex SHA-256 of the verified-boot key:
\texttt{fingerprintsNonStock} for the self-signed state and
\texttt{fingerprintsStock} for the stock state, each with a StrongBox
variant. The GrapheneOS release keys are hardcoded in the self-signed
table \citep{grapheneos-auditor-source}:

\begin{verbatim}
private static final ImmutableMap<String, DeviceInfo> fingerprintsNonStock = ImmutableMap
        .<String, DeviceInfo>builder()
        // GrapheneOS
        .put("F0A890375D1405E62EBFD87E8D3F475F948EF031BBF9DDD516D5F600A23677E8",
                new DeviceInfo(R.string.device_pixel_6, 100, 100, false, R.string.os_graphene))
        ...
\end{verbatim}

Verification looks the attested key up in the table for the reported
boot state and rejects any key that is absent, so a build signed with a
key not compiled in -\/- including an owner\textquotesingle s
self-signed build of otherwise unmodified GrapheneOS -\/- is refused
rather than attested:

\begin{verbatim}
final RootOfTrust.VerifiedBootState verifiedBootState = rootOfTrust.verifiedBootState;
final String verifiedBootKey = BaseEncoding.base16().encode(rootOfTrust.verifiedBootKey);
final DeviceInfo device;
if (verifiedBootState == RootOfTrust.VerifiedBootState.SELF_SIGNED) {
    if (attestationSecurityLevelEnum == ParsedAttestationRecord.SecurityLevel.STRONG_BOX) {
        device = fingerprintsStrongBoxNonStock.get(verifiedBootKey);
    } else {
        device = fingerprintsNonStock.get(verifiedBootKey);
    }
} else if (verifiedBootState == RootOfTrust.VerifiedBootState.VERIFIED) {
    if (attestationSecurityLevelEnum == ParsedAttestationRecord.SecurityLevel.STRONG_BOX) {
        device = fingerprintsStrongBoxStock.get(verifiedBootKey);
    } else {
        device = fingerprintsStock.get(verifiedBootKey);
    }
} else {
    throw new GeneralSecurityException("verified boot state is not verified or self signed");
}

if (device == null) {
    throw new GeneralSecurityException("invalid verified boot key fingerprint: " + verifiedBootKey);
}
\end{verbatim}

At the time of writing, GrapheneOS publishes no user-facing procedure
for adding or replacing an accepted verified-boot key with an
owner\textquotesingle s own -\/- its device-support tutorial covers only
submitting a locked-stock sample for GrapheneOS itself to add a new
factory device, not an owner or alternative-OS key
\citep{grapheneos-auditor}, \citep{grapheneos-auditor-tutorial}. It
states only that an alternative operating system\textquotesingle s
verified-boot key must be "included in the Auditor app and Attestation
Server." Recognizing an owner-signed build therefore means editing these
hardcoded tables in the source and rebuilding the Auditor app: local
second-device verification needs only that
\citep{grapheneos-auditor-tutorial}, while scheduled remote verification
\citep{grapheneos-auditor-tutorial} additionally requires rebuilding the
app to point at a different server -\/- the released Auditor pins
attestation.app and rejects any other \citep{grapheneos-auditor-domain},
\citep{grapheneos-auditor-issue140} -\/- and self-hosting a
correspondingly keyed attestation server
\citep{grapheneos-attestation-server}.

Two observations from practice bear this out. An owner who locks the
bootloader to their own AVB key and tries to enable remote verification
against the official service receives an HTTP 400 response
\citep{grapheneos-attestation-400}. And where an owner-rooted key is
accepted at all, it is by forking both components: DiamaneOS, a
downstream build, ships its own key allowlists, externalizes the trusted
keys to a replaceable configuration, and adds a second trust profile
that accepts a device reporting self-signed verified boot whose key
chains to no manufacturer root -\/- generalized owner-key attestation,
but gated by that project\textquotesingle s own release list and the
account holder\textquotesingle s approval rather than offered as an
end-user facility \citep{diamaneos-attestation}; the robotnix build
system similarly patches the Auditor app and attestation server to use
owner-created keys, automating the owner-signed OS together with a
matching app and server \citep{robotnix-attestation}.

Source: GrapheneOS Auditor,
\texttt{app/src/main/java/app/attestation/auditor/AttestationProtocol.java},
at commit d5d326453e6e3b328d7a11198d9f99174ac7f4f3
(https://github.com/GrapheneOS/Auditor). Copyright (c) 2018-2025
GrapheneOS, MIT License. These excerpts are reproduced for scholarly
commentary; they remain under their original copyright and MIT license,
not this paper\textquotesingle s license.

  \bibliography{refs}

\end{document}
